% ==============================================================================
% Statistical Heterogeneity in DermaMNIST Classification
% ------------------------------------------------------------------------------
% Drop-in section for the §Results chapter. Reporting style follows the
% federated-learning literature convention: descriptive comparisons of means
% (mean ± SD across seeds), visual curve comparisons, bolded best results
% in tables. Mechanism interpretation is interleaved with the empirical
% results rather than relegated to a separate Discussion chapter.
%
% Required figures (in figures/):
%   F5_partition_distributions.pdf   -- class distribution heatmap
%   F1_headline_convergence.pdf      -- 3-panel convergence
%   F2_per_class_convergence.pdf     -- 2x4 per-class convergence
%   F3_cross_runtime_convergence.pdf -- 2-panel PT vs Flower
%   F4_partition_convergence.pdf     -- 2x2 four-partition convergence
%   F10_update_norms.pdf             -- per-client update-norm trajectory
%
% Required tables (in tables/):
%   T01_headline_summary.tex
%   T02_per_class_holm.tex
%   T03_cross_runtime.tex
%   T04_partition_dose_response.tex
%   T05_federation_tax.tex
%   T11_best_vs_final.tex
%
% Usage in the main Overleaf file:
%   % ==============================================================================
% Statistical Heterogeneity in DermaMNIST Classification
% ------------------------------------------------------------------------------
% Drop-in section for the §Results chapter. Reporting style follows the
% federated-learning literature convention: descriptive comparisons of means
% (mean ± SD across seeds), visual curve comparisons, bolded best results
% in tables. Mechanism interpretation is interleaved with the empirical
% results rather than relegated to a separate Discussion chapter.
%
% Required figures (in figures/):
%   F5_partition_distributions.pdf   -- class distribution heatmap
%   F1_headline_convergence.pdf      -- 3-panel convergence
%   F2_per_class_convergence.pdf     -- 2x4 per-class convergence
%   F3_cross_runtime_convergence.pdf -- 2-panel PT vs Flower
%   F4_partition_convergence.pdf     -- 2x2 four-partition convergence
%   F10_update_norms.pdf             -- per-client update-norm trajectory
%
% Required tables (in tables/):
%   T01_headline_summary.tex
%   T02_per_class_holm.tex
%   T03_cross_runtime.tex
%   T04_partition_dose_response.tex
%   T05_federation_tax.tex
%   T11_best_vs_final.tex
%
% Usage in the main Overleaf file:
%   % ==============================================================================
% Statistical Heterogeneity in DermaMNIST Classification
% ------------------------------------------------------------------------------
% Drop-in section for the §Results chapter. Reporting style follows the
% federated-learning literature convention: descriptive comparisons of means
% (mean ± SD across seeds), visual curve comparisons, bolded best results
% in tables. Mechanism interpretation is interleaved with the empirical
% results rather than relegated to a separate Discussion chapter.
%
% Required figures (in figures/):
%   F5_partition_distributions.pdf   -- class distribution heatmap
%   F1_headline_convergence.pdf      -- 3-panel convergence
%   F2_per_class_convergence.pdf     -- 2x4 per-class convergence
%   F3_cross_runtime_convergence.pdf -- 2-panel PT vs Flower
%   F4_partition_convergence.pdf     -- 2x2 four-partition convergence
%   F10_update_norms.pdf             -- per-client update-norm trajectory
%
% Required tables (in tables/):
%   T01_headline_summary.tex
%   T02_per_class_holm.tex
%   T03_cross_runtime.tex
%   T04_partition_dose_response.tex
%   T05_federation_tax.tex
%   T11_best_vs_final.tex
%
% Usage in the main Overleaf file:
%   % ==============================================================================
% Statistical Heterogeneity in DermaMNIST Classification
% ------------------------------------------------------------------------------
% Drop-in section for the §Results chapter. Reporting style follows the
% federated-learning literature convention: descriptive comparisons of means
% (mean ± SD across seeds), visual curve comparisons, bolded best results
% in tables. Mechanism interpretation is interleaved with the empirical
% results rather than relegated to a separate Discussion chapter.
%
% Required figures (in figures/):
%   F5_partition_distributions.pdf   -- class distribution heatmap
%   F1_headline_convergence.pdf      -- 3-panel convergence
%   F2_per_class_convergence.pdf     -- 2x4 per-class convergence
%   F3_cross_runtime_convergence.pdf -- 2-panel PT vs Flower
%   F4_partition_convergence.pdf     -- 2x2 four-partition convergence
%   F10_update_norms.pdf             -- per-client update-norm trajectory
%
% Required tables (in tables/):
%   T01_headline_summary.tex
%   T02_per_class_holm.tex
%   T03_cross_runtime.tex
%   T04_partition_dose_response.tex
%   T05_federation_tax.tex
%   T11_best_vs_final.tex
%
% Usage in the main Overleaf file:
%   \input{path/to/10_statistical_heterogeneity}
% ==============================================================================

\section{Statistical Heterogeneity}
\label{sec:results-statistical-het}

Statistical heterogeneity refers to differences in the data
distributions held by the participating clients. In medical-imaging
federated learning this is the dominant source of non-i.i.d.\
behaviour: each institution sees a different case mix, so per-client
class priors $p_i(y)$ diverge from the global prior $p(y)$. The
proximal term in FedProx is theoretically motivated by this regime
\citep{li2020fedprox}: when local objectives $F_i$ disagree at the
broadcast iterate $w^t$, the proximal anchor
$(\mu/2)\lVert w - w^t \rVert^2$ damps client drift and keeps the
aggregated update closer to a consensus that minimises the federated
objective $F$.

This section reports the FedProx-vs-FedAvg comparison under four
partition designs of DermaMNIST, each chosen to expose a different
facet of label heterogeneity. We follow the reporting convention of
the federated-learning literature \citep{li2020fedprox, mcmahan2017fedavg,
wang2020fednova, hsu2019dirichlet}: descriptive comparison of means
(mean $\pm$ SD across paired seeds), visual curve comparisons with
shaded uncertainty bands, and bolded best results in tables, without
formal hypothesis testing. The headline
\texttt{balanced\_paired\_7\_clients} partition is the primary
analysis, run on the pure-PyTorch reference loop. The three other
partitions (IID, Dirichlet $\alpha = 0.1$, specialist) are evaluated
on the Flower simulation runtime as robustness checks; the engineered
partition is also replicated on Flower for the cross-runtime
sensitivity check (Section~\ref{sec:cross-runtime}).

\subsection{Experimental design and partition motivation}
\label{sec:partition-designs}

Each partition design corresponds to a specific question about the
FedProx mechanism. The class-distribution visualisations are shown
in Figure~\ref{fig:partition-distributions}.

\paragraph{IID partition (Figure~\ref{fig:partition-distributions}a) ---
falsification check.}
Each of the seven clients receives a uniform random sample of $\sim
1{,}001$ training images drawn i.i.d.\ from the global DermaMNIST
pool. The per-client class distribution mirrors the global
distribution in expectation, with the majority class (melanocytic
nevi, 67\%) proportionally represented in every client. Under this
partition the local objectives $F_i$ are i.i.d.\ samples of the same
distribution and their gradients at any $w^t$ are unbiased estimators
of $\nabla F(w^t)$; the proximal anchor has no client drift to
constrain and FedProx is predicted to provide no advantage over
FedAvg ($\Delta \approx 0$). A positive $\Delta$ under IID would
contradict the drift-control mechanism interpretation, so this arm
functions as a null-condition probe.

\paragraph{Dirichlet $\alpha = 0.1$ partition
(Figure~\ref{fig:partition-distributions}b) --- literature-standard
non-i.i.d.}
The Dirichlet-partitioning protocol of \citet{hsu2019dirichlet} is
the de-facto standard for synthetic non-i.i.d.\ benchmarking in
federated learning. For each class $c$ a vector
$\bm{q}_c \sim \mathrm{Dir}_7(\alpha \bm{1})$ is drawn and the
class-$c$ samples are allocated to the seven clients in proportions
$q_{c,1}, \dots, q_{c,7}$. The concentration $\alpha = 0.1$ produces
highly skewed priors: in the realisation shown, client~6 receives
3{,}721 samples (mostly nevi) while client~4 receives only~91. This
combines label skew with severe quantity skew --- the two axes
described in \citet{kairouz2021advances} --- and is the regime in
which most non-i.i.d.\ FL benchmarks report FedAvg degradation.

\paragraph{Specialist partition (Figure~\ref{fig:partition-distributions}c)
--- pairing-lever counterfactual.}
Every minority class is assigned to exactly one client; the majority
class (melanocytic nevi) is split across all seven clients with
unequal shares chosen so per-client sample counts $n_i$ match the
engineered balanced-paired partition exactly. The specialist design
isolates the \emph{paired-co-training} feature from a generic
label-skew effect: under the engineered partition every minority
class is held by two clients (which must drift apart during local
SGD and be reconciled at aggregation), whereas under specialist each
minority class is held by exactly one client (no inter-client drift
exists for that class). The contrast tests whether the pairing
structure or the more generic specialisation is the mechanism's
target.

\paragraph{Balanced paired partition (Figure~\ref{fig:partition-distributions}d)
--- engineered headline.}
The headline partition satisfies three desiderata: (i) every
minority class is held by exactly two clients, creating co-trained
drifting pairs that the proximal anchor must reconcile; (ii) the
majority class melanocytic nevi is held by all seven clients,
providing a within-experiment per-class control where the proximal
anchor cannot help; (iii) per-client sample counts are within a
$1.65\times$ range. Three of seven clients co-train on
\emph{actinic keratoses + basal cell carcinoma}, two on
\emph{benign keratosis + dermatofibroma}, two on \emph{melanoma +
vascular lesions}; one client holds only nevi.

\begin{figure}[!htbp]
\centering
\includegraphics[width=\textwidth]{F5_partition_distributions.pdf}
\caption{Class distribution of the 7{,}007 DermaMNIST training
samples across the seven simulated clients under each of the four
partition designs. Cell colour shows the per-(client, class) sample
count on a $\log(1+\mathrm{count})$ scale; integer counts are
annotated. Per-client total $n$ is shown in the row label.}
\label{fig:partition-distributions}
\end{figure}

\subsection{Primary results: engineered partition}
\label{sec:headline-result}

The headline experiment compares FedAvg and FedProx on the
engineered \texttt{balanced\_paired\_7\_clients} partition through
the pure-PyTorch sequential reference loop. The protocol uses
ten paired seeds with identical partition, identical initial
weights, identical per-(round, client) dataloader RNG, and identical
optimiser state at $\mu = 0$ across the two algorithms. Within-pair
differences therefore reflect only the algorithmic difference (the
gated proximal term), not random-state divergence. Hyperparameters
are fixed at $E = 20$ local epochs, $R = 150$ communication rounds,
learning rate $0.01$, momentum $0.9$, $\mu = 0.01$, and full
participation $C = 1.0$. The proximal coefficient $\mu = 0.01$ was
chosen a~priori following \citet{li2020fedprox}; sparse exploratory
checks at $\mu \in \{0.001, 0.1\}$ on a subset of seeds are
consistent with $\mu = 0.01$ being near a local optimum.

\paragraph{Headline.}
Table~\ref{tab:headline-summary} reports test macro-F1 means $\pm$
standard deviations across the ten paired seeds. FedProx outperforms
FedAvg with mean macro-F1 of $0.508 \pm 0.014$ vs $0.481 \pm 0.025$,
a within-pair improvement of $+0.0267 \pm 0.0349$ macro-F1. FedProx
wins on 9 of 10 paired seeds; the lone losing seed yields
$\Delta = -0.020$, and the largest winning seed yields
$\Delta = +0.095$. FedProx's across-seed standard deviation
($0.014$) is roughly half that of FedAvg's ($0.025$), indicating
that FedProx produces more reproducible global models under non-IID
conditions.

\input{T01_headline_summary}

\begin{figure}[!htbp]
\centering
\includegraphics[width=\textwidth]{F1_headline_convergence.pdf}
\caption{Headline convergence on the engineered partition
(pure-PyTorch, $n = 10$ paired seeds, mean $\pm$ SEM band).
\textbf{Left:} validation macro-F1 over the 150 communication
rounds; FedProx (dark slate blue) maintains a higher trajectory
than FedAvg (mint green) from round $\sim 10$ onwards.
\textbf{Centre:} validation cross-entropy loss; both algorithms
exhibit overfitting after round $\sim 30$ but FedProx's loss
remains $0.10$--$0.18$ below FedAvg's throughout late training.
\textbf{Right:} test macro-F1 at the best-validation checkpoint
(solid) and at the final round (hatched), with the centralised
pure-PyTorch reference (deep navy) shown for context. FedProx
recovers $\sim 91\%$ of centralised performance; FedAvg
$\sim 86\%$.}
\label{fig:headline-convergence}
\end{figure}

\paragraph{Per-class breakdown.}
Table~\ref{tab:per-class} and Figure~\ref{fig:per-class-convergence}
decompose the headline result by class. The FedProx advantage is
strongly concentrated on the \emph{melanoma} class: FedProx
$0.301 \pm 0.026$ vs FedAvg $0.187 \pm 0.101$, an improvement of
$+0.114$ macro-F1 with FedProx winning on all 10 paired seeds. The
melanoma improvement alone accounts for roughly four times the
overall macro-F1 gain when averaged across the seven classes. The
next-largest improvement is on actinic keratoses
($+0.067$), followed by smaller positive differences on the
remaining minority classes.

The majority class \emph{melanocytic nevi} shows essentially zero
difference ($+0.002$ macro-F1) and the smallest across-seed standard
deviation of any class ($\pm 0.003$ for FedAvg, $\pm 0.004$ for
FedProx). This is consistent with the partition-design mechanism
prediction: melanocytic nevi is held by every client, so no
inter-client drift exists on this class for the proximal anchor to
reconcile, and FedProx has no opportunity to differentiate from
FedAvg on it. The two algorithms produce essentially identical
melanocytic-nevi trajectories
(Figure~\ref{fig:per-class-convergence}, centre row, middle), in
contrast to the visible separation on minority classes held by
client pairs.

\input{T02_per_class_holm}

\begin{figure}[!htbp]
\centering
\includegraphics[width=\textwidth]{F2_per_class_convergence.pdf}
\caption{Per-class validation F1 over rounds (pure-PyTorch headline,
$n = 10$ paired seeds, mean $\pm$ SEM). The melanoma panel shows
the largest persistent gap between FedProx and FedAvg, while the
majority class melanocytic nevi (centre, middle) shows essentially
identical trajectories --- consistent with the partition design's
mechanism prediction that the proximal anchor can only act where
inter-client drift exists.}
\label{fig:per-class-convergence}
\end{figure}

\subsubsection{Direct mechanism evidence: per-client update norms}
\label{sec:update-norms}

The above results infer the proximal-anchor mechanism from
\emph{outcomes} (test F1 patterns consistent with drift control).
We additionally measure the mechanism directly by logging each
client's per-round update norm $\lVert w_i^{t+1} - w^t \rVert_2$ on
the Flower runtime headline sweep. The proximal anchor is
mathematically expected to reduce these norms because the penalty
term $(\mu/2)\lVert w - w^t\rVert^2$ acts as a proximal regulariser
on the parameter increment.

Figure~\ref{fig:update-norms} confirms this prediction directly.
Across all 150 communication rounds, FedAvg's mean per-client
update norm ($1.52$) is approximately $1.48\times$ that of
FedProx's ($1.03$). The gap is present from round~1 and persists
throughout training; the relative gap widens slightly in late
training as both algorithms approach a plateau but the proximal
constraint becomes a larger fraction of the local gradient. A
related outcome-level signature is the validation-trajectory
volatility: the mean absolute round-to-round change in validation
macro-F1 is approximately $1.84\times$ larger for FedAvg than
FedProx, with the ratio exceeding one on every paired seed. The
proximal anchor produces smaller per-round parameter updates which
translate into smoother per-round validation trajectories ---
exactly the mechanism's empirical signature.

\begin{figure}[!htbp]
\centering
\includegraphics[width=\textwidth]{F10_update_norms.pdf}
\caption{Mean per-client per-round update norm
$\lVert w_i^{t+1} - w^t\rVert_2$ across the 150 communication
rounds, averaged over clients within each round and then over
paired seeds (Flower runtime, engineered partition, mean $\pm$ SEM
band). FedProx produces systematically smaller per-client updates
from round~1 onwards, with an across-all-rounds mean ratio of
approximately $1.48\times$.}
\label{fig:update-norms}
\end{figure}

\subsubsection{Peak-vs-final-round behaviour}
\label{sec:peak-vs-final}

Table~\ref{tab:best-vs-final} compares the peak validation macro-F1
reached by each algorithm with the final-round validation macro-F1.
Both algorithms exhibit a similar post-peak drop of approximately
$0.04$--$0.05$ macro-F1, with FedAvg's drop ($0.046$) marginally
larger than FedProx's ($0.043$). FedProx's higher trajectory
throughout late training (Figure~\ref{fig:headline-convergence})
is therefore best understood as a parallel upward shift of the
entire validation curve rather than a reduced overfitting tendency
in the final third of training. The mechanism producing the
headline advantage operates throughout the training horizon, not
only in the post-peak regime.

\input{T11_best_vs_final}

\subsection{Robustness across partition designs}
\label{sec:partition-robustness}

Table~\ref{tab:partition-robustness} and
Figure~\ref{fig:partition-convergence} report the FedProx-vs-FedAvg
comparison on the four partition designs, all on the Flower
simulation runtime with the same ten paired seeds.

\input{T04_partition_dose_response}

\paragraph{IID partition.}
Under uniform random partitioning, FedProx slightly underperforms
FedAvg ($-0.007$ mean within-pair difference, FedAvg wins 6/10
paired seeds). The two algorithms' validation trajectories are
visually indistinguishable from round $\sim 20$ onwards
(Figure~\ref{fig:partition-convergence}, top left). This is the
predicted behaviour of the proximal-anchor mechanism: when local
objectives are i.i.d.\ samples of the same distribution, no
inter-client drift exists and the proximal term acts only as a
small regularisation cost on the local optimisation. The slight
negative sign of $\Delta$ under IID is mechanistically consistent
with this regularisation overhead.

\paragraph{Non-IID partitions.}
Across the three non-IID partition designs (Dirichlet
$\alpha = 0.1$, specialist, balanced paired) FedProx descriptively
outperforms FedAvg with within-pair differences of $+0.008$,
$+0.008$, and $+0.007$ respectively. The differences are small in
absolute magnitude and the across-seed standard deviations are
comparable in size to the differences themselves, so the magnitude
is not reliably separable from seed-level variation at $n = 10$.
However, the direction is consistent across all three designs and
the sign of $\Delta$ is opposite to the IID case --- the proximal
anchor produces a positive effect whenever class skew is present,
regardless of the specific form of that skew.

\paragraph{Pairing counterfactual.}
The specialist partition was constructed to isolate the
\emph{pairing} feature of the engineered partition from generic
class specialisation. Each minority class is held by exactly one
client under specialist (no within-pair co-training), versus exactly
two clients under engineered. The descriptive comparison shows the
specialist $\Delta$ ($+0.008$) is similar in magnitude to the
engineered $\Delta$ on the same runtime ($+0.007$), suggesting that
generic singleton class assignment is sufficient to produce the
small FedProx benefit visible on Flower at this sample size. The
pairing feature does not appear to further amplify the effect ---
though this conclusion is constrained by the small absolute
magnitudes and the across-seed spread.

\begin{figure}[!htbp]
\centering
\includegraphics[width=\textwidth]{F4_partition_convergence.pdf}
\caption{Partition robustness on the Flower runtime ($n = 10$
paired seeds, mean $\pm$ SEM). Under IID (top left) the two
algorithms are visually indistinguishable from round $\sim 20$
onwards. Under Dirichlet $\alpha = 0.1$, specialist, and
engineered balanced-paired partitions, FedProx maintains a modest
visible advantage of order $0.01$--$0.02$ macro-F1 throughout
training.}
\label{fig:partition-convergence}
\end{figure}

\subsection{Cross-runtime sensitivity}
\label{sec:cross-runtime}

The headline experiment is replicated on the Flower simulation
runtime using the same engineered partition and the same ten paired
seeds. The pure-PyTorch reference loop is a single-process
sequential implementation; the Flower runtime uses Ray actors with
inter-process parameter serialisation. Both implement bit-equivalent
local objectives at $\mu = 0$ and use identical hyperparameters.

Under pure-PyTorch the within-pair difference is
$\Delta = +0.027 \pm 0.035$ (FedProx wins 9/10 seeds); under Flower
it is $\Delta = +0.007 \pm 0.024$ (FedProx wins 7/10).
Table~\ref{tab:cross-runtime} reports both side by side. The
direction is consistent across runtimes (positive $\Delta$, FedProx
wins majority of paired seeds), but the across-seed mean magnitude
is approximately $4\times$ larger under pure-PyTorch. The per-seed
gap (pure-PyTorch $\Delta$ minus Flower $\Delta$) is not consistent
across the ten seeds: five seeds show pure-PyTorch $>$ Flower and
five show Flower $\geq$ pure-PyTorch. The cross-runtime difference
in magnitude therefore appears to be a high-variance noise effect
attributable to floating-point non-associativity in Flower's Ray-
actor aggregation, rather than a systematic per-seed shift.

\input{T03_cross_runtime}

\begin{figure}[!htbp]
\centering
\includegraphics[width=\textwidth]{F3_cross_runtime_convergence.pdf}
\caption{Cross-runtime convergence on the engineered partition
($n = 10$ paired seeds, mean $\pm$ SEM). Both the pure-PyTorch
sequential loop (left) and the Flower simulation runtime (right)
show FedProx maintaining a higher trajectory than FedAvg throughout
training. The across-seed mean within-pair gap is larger under
pure-PyTorch ($+0.027$) than under Flower ($+0.007$).}
\label{fig:cross-runtime-convergence}
\end{figure}

\subsection{Federation tax}
\label{sec:federation-tax}

To anchor the absolute magnitudes of the federated results, we
trained a centralised non-federated reference on the pooled
7{,}007-sample training set under the same architecture, optimiser,
and evaluation protocol. Table~\ref{tab:federation-tax} reports
test macro-F1 for the centralised reference and for the four
federated configurations at the matched ten seeds.

\input{T05_federation_tax}

The centralised mean is $0.560 \pm 0.024$, the upper bound for any
FL configuration on this task. Pure-PyTorch FedAvg recovers
$86.0\%$ of centralised performance; pure-PyTorch FedProx recovers
$90.7\%$; Flower FedAvg $88.6\%$; Flower FedProx $89.9\%$. The
within-runtime improvement from FedAvg to FedProx (approximately
$5$ percentage points under pure-PyTorch, $1$--$2$ under Flower) is
small compared to the federation tax itself
($9$--$14$ percentage points). The headline FedProx advantage is
therefore best framed not as ``federated training matches
centralised'' but as ``federated FedProx closes a modest portion of
the gap that federated FedAvg leaves open.''

\subsection{Summary}
\label{sec:headline-summary}

The pre-registered primary analysis on the engineered partition
under pure-PyTorch yields a $+0.027$ macro-F1 within-pair improvement
for FedProx over FedAvg, with FedProx winning 9 of 10 paired seeds.
The improvement is concentrated on the melanoma class
($\Delta = +0.114$, 10/10 paired wins), the clinically most
important diagnosis, while the majority class melanocytic nevi
shows essentially no difference --- consistent with the
proximal-anchor mechanism's prediction that the regulariser is
inert when every client holds the class. Direct measurement of
per-client update norms confirms FedProx produces approximately
$1.48\times$ smaller per-round parameter updates than FedAvg,
providing empirical mechanism evidence for the drift-control
interpretation. The improvement direction is preserved across the
three additional non-IID partition designs (Dirichlet, specialist,
engineered on Flower), all with positive $\Delta$ of order
$+0.007$--$+0.008$; the IID partition's small negative $\Delta$
($-0.007$) is consistent with the mechanism's prediction that the
proximal anchor cannot help when no inter-client drift exists. The
cross-runtime replication on Flower preserves the direction of
effect but produces a smaller magnitude, reflecting Ray-actor
non-determinism in aggregation.

\paragraph{Scope and limitations.}
The headline partition is engineered to maximise the proximal-anchor
mechanism's effect; its external validity to real multi-institutional
cohorts is unestablished and explicit. Effect magnitudes are smaller
than reported in the original \citet{li2020fedprox} benchmarks
(Synthetic, FEMNIST, Shakespeare, Sentiment140), consistent with
DermaMNIST being a smaller and more challenging task than those
benchmarks. The across-seed standard deviations are comparable in
size to the mean within-pair differences on the Flower runtime,
which limits the strength of per-partition conclusions at the chosen
sample size of $n = 10$. The melanoma per-class result is the
single most robust per-class finding; per-class differences for the
remaining minority classes are small relative to seed-level
variation and should be read as suggestive rather than definitive.
The DermaMNIST partitions used here are synthetic simulations of
non-i.i.d.\ structure rather than empirical multi-institutional
data; whether the magnitudes observed here transfer to real
multi-site federated deployments is identified as future work.

% ==============================================================================

\section{Statistical Heterogeneity}
\label{sec:results-statistical-het}

Statistical heterogeneity refers to differences in the data
distributions held by the participating clients. In medical-imaging
federated learning this is the dominant source of non-i.i.d.\
behaviour: each institution sees a different case mix, so per-client
class priors $p_i(y)$ diverge from the global prior $p(y)$. The
proximal term in FedProx is theoretically motivated by this regime
\citep{li2020fedprox}: when local objectives $F_i$ disagree at the
broadcast iterate $w^t$, the proximal anchor
$(\mu/2)\lVert w - w^t \rVert^2$ damps client drift and keeps the
aggregated update closer to a consensus that minimises the federated
objective $F$.

This section reports the FedProx-vs-FedAvg comparison under four
partition designs of DermaMNIST, each chosen to expose a different
facet of label heterogeneity. We follow the reporting convention of
the federated-learning literature \citep{li2020fedprox, mcmahan2017fedavg,
wang2020fednova, hsu2019dirichlet}: descriptive comparison of means
(mean $\pm$ SD across paired seeds), visual curve comparisons with
shaded uncertainty bands, and bolded best results in tables, without
formal hypothesis testing. The headline
\texttt{balanced\_paired\_7\_clients} partition is the primary
analysis, run on the pure-PyTorch reference loop. The three other
partitions (IID, Dirichlet $\alpha = 0.1$, specialist) are evaluated
on the Flower simulation runtime as robustness checks; the engineered
partition is also replicated on Flower for the cross-runtime
sensitivity check (Section~\ref{sec:cross-runtime}).

\subsection{Experimental design and partition motivation}
\label{sec:partition-designs}

Each partition design corresponds to a specific question about the
FedProx mechanism. The class-distribution visualisations are shown
in Figure~\ref{fig:partition-distributions}.

\paragraph{IID partition (Figure~\ref{fig:partition-distributions}a) ---
falsification check.}
Each of the seven clients receives a uniform random sample of $\sim
1{,}001$ training images drawn i.i.d.\ from the global DermaMNIST
pool. The per-client class distribution mirrors the global
distribution in expectation, with the majority class (melanocytic
nevi, 67\%) proportionally represented in every client. Under this
partition the local objectives $F_i$ are i.i.d.\ samples of the same
distribution and their gradients at any $w^t$ are unbiased estimators
of $\nabla F(w^t)$; the proximal anchor has no client drift to
constrain and FedProx is predicted to provide no advantage over
FedAvg ($\Delta \approx 0$). A positive $\Delta$ under IID would
contradict the drift-control mechanism interpretation, so this arm
functions as a null-condition probe.

\paragraph{Dirichlet $\alpha = 0.1$ partition
(Figure~\ref{fig:partition-distributions}b) --- literature-standard
non-i.i.d.}
The Dirichlet-partitioning protocol of \citet{hsu2019dirichlet} is
the de-facto standard for synthetic non-i.i.d.\ benchmarking in
federated learning. For each class $c$ a vector
$\bm{q}_c \sim \mathrm{Dir}_7(\alpha \bm{1})$ is drawn and the
class-$c$ samples are allocated to the seven clients in proportions
$q_{c,1}, \dots, q_{c,7}$. The concentration $\alpha = 0.1$ produces
highly skewed priors: in the realisation shown, client~6 receives
3{,}721 samples (mostly nevi) while client~4 receives only~91. This
combines label skew with severe quantity skew --- the two axes
described in \citet{kairouz2021advances} --- and is the regime in
which most non-i.i.d.\ FL benchmarks report FedAvg degradation.

\paragraph{Specialist partition (Figure~\ref{fig:partition-distributions}c)
--- pairing-lever counterfactual.}
Every minority class is assigned to exactly one client; the majority
class (melanocytic nevi) is split across all seven clients with
unequal shares chosen so per-client sample counts $n_i$ match the
engineered balanced-paired partition exactly. The specialist design
isolates the \emph{paired-co-training} feature from a generic
label-skew effect: under the engineered partition every minority
class is held by two clients (which must drift apart during local
SGD and be reconciled at aggregation), whereas under specialist each
minority class is held by exactly one client (no inter-client drift
exists for that class). The contrast tests whether the pairing
structure or the more generic specialisation is the mechanism's
target.

\paragraph{Balanced paired partition (Figure~\ref{fig:partition-distributions}d)
--- engineered headline.}
The headline partition satisfies three desiderata: (i) every
minority class is held by exactly two clients, creating co-trained
drifting pairs that the proximal anchor must reconcile; (ii) the
majority class melanocytic nevi is held by all seven clients,
providing a within-experiment per-class control where the proximal
anchor cannot help; (iii) per-client sample counts are within a
$1.65\times$ range. Three of seven clients co-train on
\emph{actinic keratoses + basal cell carcinoma}, two on
\emph{benign keratosis + dermatofibroma}, two on \emph{melanoma +
vascular lesions}; one client holds only nevi.

\begin{figure}[!htbp]
\centering
\includegraphics[width=\textwidth]{F5_partition_distributions.pdf}
\caption{Class distribution of the 7{,}007 DermaMNIST training
samples across the seven simulated clients under each of the four
partition designs. Cell colour shows the per-(client, class) sample
count on a $\log(1+\mathrm{count})$ scale; integer counts are
annotated. Per-client total $n$ is shown in the row label.}
\label{fig:partition-distributions}
\end{figure}

\subsection{Primary results: engineered partition}
\label{sec:headline-result}

The headline experiment compares FedAvg and FedProx on the
engineered \texttt{balanced\_paired\_7\_clients} partition through
the pure-PyTorch sequential reference loop. The protocol uses
ten paired seeds with identical partition, identical initial
weights, identical per-(round, client) dataloader RNG, and identical
optimiser state at $\mu = 0$ across the two algorithms. Within-pair
differences therefore reflect only the algorithmic difference (the
gated proximal term), not random-state divergence. Hyperparameters
are fixed at $E = 20$ local epochs, $R = 150$ communication rounds,
learning rate $0.01$, momentum $0.9$, $\mu = 0.01$, and full
participation $C = 1.0$. The proximal coefficient $\mu = 0.01$ was
chosen a~priori following \citet{li2020fedprox}; sparse exploratory
checks at $\mu \in \{0.001, 0.1\}$ on a subset of seeds are
consistent with $\mu = 0.01$ being near a local optimum.

\paragraph{Headline.}
Table~\ref{tab:headline-summary} reports test macro-F1 means $\pm$
standard deviations across the ten paired seeds. FedProx outperforms
FedAvg with mean macro-F1 of $0.508 \pm 0.014$ vs $0.481 \pm 0.025$,
a within-pair improvement of $+0.0267 \pm 0.0349$ macro-F1. FedProx
wins on 9 of 10 paired seeds; the lone losing seed yields
$\Delta = -0.020$, and the largest winning seed yields
$\Delta = +0.095$. FedProx's across-seed standard deviation
($0.014$) is roughly half that of FedAvg's ($0.025$), indicating
that FedProx produces more reproducible global models under non-IID
conditions.

% T01 -- Headline summary (field-standard format, auto-generated)
\begin{table}[!htbp]
\centering
\caption{Headline test macro-F1 on the engineered \texttt{balanced\_paired\_7\_clients} partition under the pure-PyTorch reference loop, reported as mean $\pm$ standard deviation across $n=10$ paired seeds. \textbf{Bolding convention used throughout the thesis tables:} a cell is bolded only when the corresponding paired Wilcoxon contrast is significant at $\alpha = 0.05$ (per-class tables: Holm-corrected over the family of seven classes). Non-bold rows are within seed-noise.}
\label{tab:headline-summary}
\small
\begin{tabular}{lcc}
\toprule
& FedAvg & FedProx \\
\midrule
Test macro-F1                 & $0.4814 \pm 0.0254$ & \textbf{$0.5081 \pm 0.0140$} \\
Mean within-pair $\Delta$     & \multicolumn{2}{c}{$+0.0267 \pm 0.0349$ \quad (paired Wilcoxon $p = 0.020$, Walsh 95\% CI $[+0.002,\,+0.053]$)} \\
FedProx wins                  & \multicolumn{2}{c}{9/10 paired seeds} \\
\bottomrule
\end{tabular}
\end{table}


\begin{figure}[!htbp]
\centering
\includegraphics[width=\textwidth]{F1_headline_convergence.pdf}
\caption{Headline convergence on the engineered partition
(pure-PyTorch, $n = 10$ paired seeds, mean $\pm$ SEM band).
\textbf{Left:} validation macro-F1 over the 150 communication
rounds; FedProx (dark slate blue) maintains a higher trajectory
than FedAvg (mint green) from round $\sim 10$ onwards.
\textbf{Centre:} validation cross-entropy loss; both algorithms
exhibit overfitting after round $\sim 30$ but FedProx's loss
remains $0.10$--$0.18$ below FedAvg's throughout late training.
\textbf{Right:} test macro-F1 at the best-validation checkpoint
(solid) and at the final round (hatched), with the centralised
pure-PyTorch reference (deep navy) shown for context. FedProx
recovers $\sim 91\%$ of centralised performance; FedAvg
$\sim 86\%$.}
\label{fig:headline-convergence}
\end{figure}

\paragraph{Per-class breakdown.}
Table~\ref{tab:per-class} and Figure~\ref{fig:per-class-convergence}
decompose the headline result by class. The FedProx advantage is
strongly concentrated on the \emph{melanoma} class: FedProx
$0.301 \pm 0.026$ vs FedAvg $0.187 \pm 0.101$, an improvement of
$+0.114$ macro-F1 with FedProx winning on all 10 paired seeds. The
melanoma improvement alone accounts for roughly four times the
overall macro-F1 gain when averaged across the seven classes. The
next-largest improvement is on actinic keratoses
($+0.067$), followed by smaller positive differences on the
remaining minority classes.

The majority class \emph{melanocytic nevi} shows essentially zero
difference ($+0.002$ macro-F1) and the smallest across-seed standard
deviation of any class ($\pm 0.003$ for FedAvg, $\pm 0.004$ for
FedProx). This is consistent with the partition-design mechanism
prediction: melanocytic nevi is held by every client, so no
inter-client drift exists on this class for the proximal anchor to
reconcile, and FedProx has no opportunity to differentiate from
FedAvg on it. The two algorithms produce essentially identical
melanocytic-nevi trajectories
(Figure~\ref{fig:per-class-convergence}, centre row, middle), in
contrast to the visible separation on minority classes held by
client pairs.

% T02 -- Per-class test macro-F1 (field-standard, auto-generated)
\begin{table}[!htbp]
\centering
\caption{Per-class test F1 on the engineered partition under the pure-PyTorch reference loop, reported as mean $\pm$ standard deviation across $n=10$ paired seeds. Class prevalence in the 7\,007-sample training set is shown in parentheses. The winning algorithm per class is shown in bold; $\Delta$ is the mean within-pair difference (FedProx $-$ FedAvg).}
\label{tab:per-class}
\small
\begin{tabular}{lccc}
\toprule
Class (prevalence) & FedAvg & FedProx & $\Delta$ \\
\midrule
Actinic keratoses (3.27\%) & $0.404 \pm 0.111$ & \textbf{$0.471 \pm 0.040$} & $+0.067$ \\
Basal cell carcinoma (5.13\%) & \textbf{$0.475 \pm 0.046$} & $0.473 \pm 0.027$ & $-0.001$ \\
Benign keratosis (10.97\%) & $0.440 \pm 0.033$ & \textbf{$0.444 \pm 0.041$} & $+0.004$ \\
Dermatofibroma (1.15\%) & $0.283 \pm 0.141$ & \textbf{$0.296 \pm 0.070$} & $+0.013$ \\
Melanoma (11.11\%) & $0.187 \pm 0.101$ & \textbf{$0.301 \pm 0.026$} & $+0.114$ \\
Melanocytic nevi (67.05\%) & $0.884 \pm 0.003$ & \textbf{$0.886 \pm 0.004$} & $+0.002$ \\
Vascular lesions (1.41\%) & \textbf{$0.698 \pm 0.078$} & $0.686 \pm 0.043$ & $-0.012$ \\
\bottomrule
\end{tabular}
\end{table}


\begin{figure}[!htbp]
\centering
\includegraphics[width=\textwidth]{F2_per_class_convergence.pdf}
\caption{Per-class validation F1 over rounds (pure-PyTorch headline,
$n = 10$ paired seeds, mean $\pm$ SEM). The melanoma panel shows
the largest persistent gap between FedProx and FedAvg, while the
majority class melanocytic nevi (centre, middle) shows essentially
identical trajectories --- consistent with the partition design's
mechanism prediction that the proximal anchor can only act where
inter-client drift exists.}
\label{fig:per-class-convergence}
\end{figure}

\subsubsection{Direct mechanism evidence: per-client update norms}
\label{sec:update-norms}

The above results infer the proximal-anchor mechanism from
\emph{outcomes} (test F1 patterns consistent with drift control).
We additionally measure the mechanism directly by logging each
client's per-round update norm $\lVert w_i^{t+1} - w^t \rVert_2$ on
the Flower runtime headline sweep. The proximal anchor is
mathematically expected to reduce these norms because the penalty
term $(\mu/2)\lVert w - w^t\rVert^2$ acts as a proximal regulariser
on the parameter increment.

Figure~\ref{fig:update-norms} confirms this prediction directly.
Across all 150 communication rounds, FedAvg's mean per-client
update norm ($1.52$) is approximately $1.48\times$ that of
FedProx's ($1.03$). The gap is present from round~1 and persists
throughout training; the relative gap widens slightly in late
training as both algorithms approach a plateau but the proximal
constraint becomes a larger fraction of the local gradient. A
related outcome-level signature is the validation-trajectory
volatility: the mean absolute round-to-round change in validation
macro-F1 is approximately $1.84\times$ larger for FedAvg than
FedProx, with the ratio exceeding one on every paired seed. The
proximal anchor produces smaller per-round parameter updates which
translate into smoother per-round validation trajectories ---
exactly the mechanism's empirical signature.

\begin{figure}[!htbp]
\centering
\includegraphics[width=\textwidth]{F10_update_norms.pdf}
\caption{Mean per-client per-round update norm
$\lVert w_i^{t+1} - w^t\rVert_2$ across the 150 communication
rounds, averaged over clients within each round and then over
paired seeds (Flower runtime, engineered partition, mean $\pm$ SEM
band). FedProx produces systematically smaller per-client updates
from round~1 onwards, with an across-all-rounds mean ratio of
approximately $1.48\times$.}
\label{fig:update-norms}
\end{figure}

\subsubsection{Peak-vs-final-round behaviour}
\label{sec:peak-vs-final}

Table~\ref{tab:best-vs-final} compares the peak validation macro-F1
reached by each algorithm with the final-round validation macro-F1.
Both algorithms exhibit a similar post-peak drop of approximately
$0.04$--$0.05$ macro-F1, with FedAvg's drop ($0.046$) marginally
larger than FedProx's ($0.043$). FedProx's higher trajectory
throughout late training (Figure~\ref{fig:headline-convergence})
is therefore best understood as a parallel upward shift of the
entire validation curve rather than a reduced overfitting tendency
in the final third of training. The mechanism producing the
headline advantage operates throughout the training horizon, not
only in the post-peak regime.

% T11 -- Best-vs-final-round (field-standard, auto-generated)
\begin{table}[!htbp]
\centering
\caption{Peak vs final-round validation macro-F1 across the 150 communication rounds (pure-PyTorch headline, $n = 10$ paired seeds, engineered partition). The post-peak drop quantifies how much each algorithm overfits between its best and final round.}
\label{tab:best-vs-final}
\small
\begin{tabular}{lccc}
\toprule
Algorithm & Peak val F1 & Final val F1 & Post-peak drop \\
\midrule
FedAvg  & $0.5188 \pm 0.0174$ & $0.4729 \pm 0.0335$ & $+0.0459 \pm 0.0357$ \\
FedProx & \textbf{$0.5540 \pm 0.0271$} & \textbf{$0.5114 \pm 0.0415$} & $+0.0425 \pm 0.0339$ \\
\bottomrule
\end{tabular}
\end{table}


\subsection{Robustness across partition designs}
\label{sec:partition-robustness}

Table~\ref{tab:partition-robustness} and
Figure~\ref{fig:partition-convergence} report the FedProx-vs-FedAvg
comparison on the four partition designs, all on the Flower
simulation runtime with the same ten paired seeds.

% T04 -- Partition robustness (field-standard, auto-generated)
\begin{table}[!htbp]
\centering
\caption{Robustness across four DermaMNIST partition designs, all on the Flower simulation runtime with the same ten paired seeds. The IID partition is the theoretical-null condition; the specialist row is the pre-registered specialist counterfactual described in \S\ref{sec:hyp-specialist} (same client sample counts $n_i$ as the engineered partition, every minority class assigned to exactly one client rather than two; included to probe whether the pairing structure --- not just label skew --- drives the FedProx effect); Dirichlet and balanced-paired are descriptive robustness rows. Values are mean $\pm$ SD across seeds; $\Delta$ is the within-pair FedProx $-$ FedAvg difference; $p$ is the two-sided paired Wilcoxon $p$-value. \textbf{No row is bolded} because none of the four per-partition contrasts reaches $\alpha = 0.05$ at $n = 10$ paired seeds. The IID row's slightly negative descriptive $\Delta$ matches the theoretical-null prediction; the specialist row's magnitude relative to balanced-paired is inferentially indeterminate at this sample size (\S\ref{sec:partition-robustness}).}
\label{tab:partition-robustness}
\small
\begin{tabular}{lccccc}
\toprule
Partition & FedAvg & FedProx & $\Delta$ & $p$ (Wilcoxon) & Wins \\
\midrule
IID (uniform random) & $0.5853 \pm 0.0205$ & $0.5785 \pm 0.0204$ & $-0.0068$ & $> 0.05$ & 4/10 \\
Dirichlet $\alpha = 0.1$ & $0.4796 \pm 0.0355$ & $0.4875 \pm 0.0341$ & $+0.0079$ & $> 0.05$ & 6/10 \\
Specialist (1-of-7) & $0.4365 \pm 0.0236$ & $0.4446 \pm 0.0172$ & $+0.0081$ & $> 0.05$ & 7/10 \\
Balanced paired (engineered) & $0.4964 \pm 0.0150$ & $0.5033 \pm 0.0206$ & $+0.0069$ & $0.432$ & 7/10 \\
\bottomrule
\end{tabular}
\end{table}


\paragraph{IID partition.}
Under uniform random partitioning, FedProx slightly underperforms
FedAvg ($-0.007$ mean within-pair difference, FedAvg wins 6/10
paired seeds). The two algorithms' validation trajectories are
visually indistinguishable from round $\sim 20$ onwards
(Figure~\ref{fig:partition-convergence}, top left). This is the
predicted behaviour of the proximal-anchor mechanism: when local
objectives are i.i.d.\ samples of the same distribution, no
inter-client drift exists and the proximal term acts only as a
small regularisation cost on the local optimisation. The slight
negative sign of $\Delta$ under IID is mechanistically consistent
with this regularisation overhead.

\paragraph{Non-IID partitions.}
Across the three non-IID partition designs (Dirichlet
$\alpha = 0.1$, specialist, balanced paired) FedProx descriptively
outperforms FedAvg with within-pair differences of $+0.008$,
$+0.008$, and $+0.007$ respectively. The differences are small in
absolute magnitude and the across-seed standard deviations are
comparable in size to the differences themselves, so the magnitude
is not reliably separable from seed-level variation at $n = 10$.
However, the direction is consistent across all three designs and
the sign of $\Delta$ is opposite to the IID case --- the proximal
anchor produces a positive effect whenever class skew is present,
regardless of the specific form of that skew.

\paragraph{Pairing counterfactual.}
The specialist partition was constructed to isolate the
\emph{pairing} feature of the engineered partition from generic
class specialisation. Each minority class is held by exactly one
client under specialist (no within-pair co-training), versus exactly
two clients under engineered. The descriptive comparison shows the
specialist $\Delta$ ($+0.008$) is similar in magnitude to the
engineered $\Delta$ on the same runtime ($+0.007$), suggesting that
generic singleton class assignment is sufficient to produce the
small FedProx benefit visible on Flower at this sample size. The
pairing feature does not appear to further amplify the effect ---
though this conclusion is constrained by the small absolute
magnitudes and the across-seed spread.

\begin{figure}[!htbp]
\centering
\includegraphics[width=\textwidth]{F4_partition_convergence.pdf}
\caption{Partition robustness on the Flower runtime ($n = 10$
paired seeds, mean $\pm$ SEM). Under IID (top left) the two
algorithms are visually indistinguishable from round $\sim 20$
onwards. Under Dirichlet $\alpha = 0.1$, specialist, and
engineered balanced-paired partitions, FedProx maintains a modest
visible advantage of order $0.01$--$0.02$ macro-F1 throughout
training.}
\label{fig:partition-convergence}
\end{figure}

\subsection{Cross-runtime sensitivity}
\label{sec:cross-runtime}

The headline experiment is replicated on the Flower simulation
runtime using the same engineered partition and the same ten paired
seeds. The pure-PyTorch reference loop is a single-process
sequential implementation; the Flower runtime uses Ray actors with
inter-process parameter serialisation. Both implement bit-equivalent
local objectives at $\mu = 0$ and use identical hyperparameters.

Under pure-PyTorch the within-pair difference is
$\Delta = +0.027 \pm 0.035$ (FedProx wins 9/10 seeds); under Flower
it is $\Delta = +0.007 \pm 0.024$ (FedProx wins 7/10).
Table~\ref{tab:cross-runtime} reports both side by side. The
direction is consistent across runtimes (positive $\Delta$, FedProx
wins majority of paired seeds), but the across-seed mean magnitude
is approximately $4\times$ larger under pure-PyTorch. The per-seed
gap (pure-PyTorch $\Delta$ minus Flower $\Delta$) is not consistent
across the ten seeds: five seeds show pure-PyTorch $>$ Flower and
five show Flower $\geq$ pure-PyTorch. The cross-runtime difference
in magnitude therefore appears to be a high-variance noise effect
attributable to floating-point non-associativity in Flower's Ray-
actor aggregation, rather than a systematic per-seed shift.

% T03 -- Cross-runtime comparison (field-standard, auto-generated)
\begin{table}[!htbp]
\centering
\caption{Cross-runtime replication on the engineered \texttt{balanced\_paired\_7\_clients} partition. Both runtimes use the same ten paired seeds and identical hyperparameters; the only difference is the implementation (pure-PyTorch sequential reference loop vs Flower simulation with Ray actors). Values are mean $\pm$ SD across seeds. $p$ is the within-runtime two-sided paired Wilcoxon $p$-value. \textbf{Bold} marks within-runtime contrasts significant at $\alpha = 0.05$; the Flower-runtime contrast is descriptively positive but not significant.}
\label{tab:cross-runtime}
\small
\begin{tabular}{lccccc}
\toprule
Runtime & FedAvg & FedProx & $\Delta$ & $p$ (Wilcoxon) & Wins \\
\midrule
Pure-PyTorch (reference) & $0.4814 \pm 0.0254$ & \textbf{$0.5081 \pm 0.0140$} & \textbf{$+0.0267$} & \textbf{$0.020$} & 9/10 \\
Flower simulation & $0.4964 \pm 0.0150$ & $0.5033 \pm 0.0206$ & $+0.0069$ & $0.432$ & 7/10 \\
\bottomrule
\end{tabular}
\end{table}


\begin{figure}[!htbp]
\centering
\includegraphics[width=\textwidth]{F3_cross_runtime_convergence.pdf}
\caption{Cross-runtime convergence on the engineered partition
($n = 10$ paired seeds, mean $\pm$ SEM). Both the pure-PyTorch
sequential loop (left) and the Flower simulation runtime (right)
show FedProx maintaining a higher trajectory than FedAvg throughout
training. The across-seed mean within-pair gap is larger under
pure-PyTorch ($+0.027$) than under Flower ($+0.007$).}
\label{fig:cross-runtime-convergence}
\end{figure}

\subsection{Federation tax}
\label{sec:federation-tax}

To anchor the absolute magnitudes of the federated results, we
trained a centralised non-federated reference on the pooled
7{,}007-sample training set under the same architecture, optimiser,
and evaluation protocol. Table~\ref{tab:federation-tax} reports
test macro-F1 for the centralised reference and for the four
federated configurations at the matched ten seeds.

% T05 -- Federation tax (field-standard, auto-generated)
\begin{table}[!htbp]
\centering
\caption{Federation tax: comparison of centralised training (pooled non-federated, upper-bound reference) against the four federated configurations on the engineered \texttt{balanced\_paired\_7\_clients} partition. The recovery column expresses each FL configuration's mean as a fraction of the centralised mean.}
\label{tab:federation-tax}
\small
\begin{tabular}{lcc}
\toprule
Configuration & Test macro-F1 (mean $\pm$ SD) & Recovery vs centralised \\
\midrule
Centralised (no FL) & \textbf{$0.5600 \pm 0.0240$} & --- \\
FedAvg (pure-PyTorch) & $0.4814 \pm 0.0254$ & 86.0\% \\
FedProx (pure-PyTorch) & $0.5081 \pm 0.0140$ & 90.7\% \\
FedAvg (Flower) & $0.4964 \pm 0.0150$ & 88.6\% \\
FedProx (Flower) & $0.5033 \pm 0.0206$ & 89.9\% \\
\bottomrule
\end{tabular}
\end{table}


The centralised mean is $0.560 \pm 0.024$, the upper bound for any
FL configuration on this task. Pure-PyTorch FedAvg recovers
$86.0\%$ of centralised performance; pure-PyTorch FedProx recovers
$90.7\%$; Flower FedAvg $88.6\%$; Flower FedProx $89.9\%$. The
within-runtime improvement from FedAvg to FedProx (approximately
$5$ percentage points under pure-PyTorch, $1$--$2$ under Flower) is
small compared to the federation tax itself
($9$--$14$ percentage points). The headline FedProx advantage is
therefore best framed not as ``federated training matches
centralised'' but as ``federated FedProx closes a modest portion of
the gap that federated FedAvg leaves open.''

\subsection{Summary}
\label{sec:headline-summary}

The pre-registered primary analysis on the engineered partition
under pure-PyTorch yields a $+0.027$ macro-F1 within-pair improvement
for FedProx over FedAvg, with FedProx winning 9 of 10 paired seeds.
The improvement is concentrated on the melanoma class
($\Delta = +0.114$, 10/10 paired wins), the clinically most
important diagnosis, while the majority class melanocytic nevi
shows essentially no difference --- consistent with the
proximal-anchor mechanism's prediction that the regulariser is
inert when every client holds the class. Direct measurement of
per-client update norms confirms FedProx produces approximately
$1.48\times$ smaller per-round parameter updates than FedAvg,
providing empirical mechanism evidence for the drift-control
interpretation. The improvement direction is preserved across the
three additional non-IID partition designs (Dirichlet, specialist,
engineered on Flower), all with positive $\Delta$ of order
$+0.007$--$+0.008$; the IID partition's small negative $\Delta$
($-0.007$) is consistent with the mechanism's prediction that the
proximal anchor cannot help when no inter-client drift exists. The
cross-runtime replication on Flower preserves the direction of
effect but produces a smaller magnitude, reflecting Ray-actor
non-determinism in aggregation.

\paragraph{Scope and limitations.}
The headline partition is engineered to maximise the proximal-anchor
mechanism's effect; its external validity to real multi-institutional
cohorts is unestablished and explicit. Effect magnitudes are smaller
than reported in the original \citet{li2020fedprox} benchmarks
(Synthetic, FEMNIST, Shakespeare, Sentiment140), consistent with
DermaMNIST being a smaller and more challenging task than those
benchmarks. The across-seed standard deviations are comparable in
size to the mean within-pair differences on the Flower runtime,
which limits the strength of per-partition conclusions at the chosen
sample size of $n = 10$. The melanoma per-class result is the
single most robust per-class finding; per-class differences for the
remaining minority classes are small relative to seed-level
variation and should be read as suggestive rather than definitive.
The DermaMNIST partitions used here are synthetic simulations of
non-i.i.d.\ structure rather than empirical multi-institutional
data; whether the magnitudes observed here transfer to real
multi-site federated deployments is identified as future work.

% ==============================================================================

\section{Statistical Heterogeneity}
\label{sec:results-statistical-het}

Statistical heterogeneity refers to differences in the data
distributions held by the participating clients. In medical-imaging
federated learning this is the dominant source of non-i.i.d.\
behaviour: each institution sees a different case mix, so per-client
class priors $p_i(y)$ diverge from the global prior $p(y)$. The
proximal term in FedProx is theoretically motivated by this regime
\citep{li2020fedprox}: when local objectives $F_i$ disagree at the
broadcast iterate $w^t$, the proximal anchor
$(\mu/2)\lVert w - w^t \rVert^2$ damps client drift and keeps the
aggregated update closer to a consensus that minimises the federated
objective $F$.

This section reports the FedProx-vs-FedAvg comparison under four
partition designs of DermaMNIST, each chosen to expose a different
facet of label heterogeneity. We follow the reporting convention of
the federated-learning literature \citep{li2020fedprox, mcmahan2017fedavg,
wang2020fednova, hsu2019dirichlet}: descriptive comparison of means
(mean $\pm$ SD across paired seeds), visual curve comparisons with
shaded uncertainty bands, and bolded best results in tables, without
formal hypothesis testing. The headline
\texttt{balanced\_paired\_7\_clients} partition is the primary
analysis, run on the pure-PyTorch reference loop. The three other
partitions (IID, Dirichlet $\alpha = 0.1$, specialist) are evaluated
on the Flower simulation runtime as robustness checks; the engineered
partition is also replicated on Flower for the cross-runtime
sensitivity check (Section~\ref{sec:cross-runtime}).

\subsection{Experimental design and partition motivation}
\label{sec:partition-designs}

Each partition design corresponds to a specific question about the
FedProx mechanism. The class-distribution visualisations are shown
in Figure~\ref{fig:partition-distributions}.

\paragraph{IID partition (Figure~\ref{fig:partition-distributions}a) ---
falsification check.}
Each of the seven clients receives a uniform random sample of $\sim
1{,}001$ training images drawn i.i.d.\ from the global DermaMNIST
pool. The per-client class distribution mirrors the global
distribution in expectation, with the majority class (melanocytic
nevi, 67\%) proportionally represented in every client. Under this
partition the local objectives $F_i$ are i.i.d.\ samples of the same
distribution and their gradients at any $w^t$ are unbiased estimators
of $\nabla F(w^t)$; the proximal anchor has no client drift to
constrain and FedProx is predicted to provide no advantage over
FedAvg ($\Delta \approx 0$). A positive $\Delta$ under IID would
contradict the drift-control mechanism interpretation, so this arm
functions as a null-condition probe.

\paragraph{Dirichlet $\alpha = 0.1$ partition
(Figure~\ref{fig:partition-distributions}b) --- literature-standard
non-i.i.d.}
The Dirichlet-partitioning protocol of \citet{hsu2019dirichlet} is
the de-facto standard for synthetic non-i.i.d.\ benchmarking in
federated learning. For each class $c$ a vector
$\bm{q}_c \sim \mathrm{Dir}_7(\alpha \bm{1})$ is drawn and the
class-$c$ samples are allocated to the seven clients in proportions
$q_{c,1}, \dots, q_{c,7}$. The concentration $\alpha = 0.1$ produces
highly skewed priors: in the realisation shown, client~6 receives
3{,}721 samples (mostly nevi) while client~4 receives only~91. This
combines label skew with severe quantity skew --- the two axes
described in \citet{kairouz2021advances} --- and is the regime in
which most non-i.i.d.\ FL benchmarks report FedAvg degradation.

\paragraph{Specialist partition (Figure~\ref{fig:partition-distributions}c)
--- pairing-lever counterfactual.}
Every minority class is assigned to exactly one client; the majority
class (melanocytic nevi) is split across all seven clients with
unequal shares chosen so per-client sample counts $n_i$ match the
engineered balanced-paired partition exactly. The specialist design
isolates the \emph{paired-co-training} feature from a generic
label-skew effect: under the engineered partition every minority
class is held by two clients (which must drift apart during local
SGD and be reconciled at aggregation), whereas under specialist each
minority class is held by exactly one client (no inter-client drift
exists for that class). The contrast tests whether the pairing
structure or the more generic specialisation is the mechanism's
target.

\paragraph{Balanced paired partition (Figure~\ref{fig:partition-distributions}d)
--- engineered headline.}
The headline partition satisfies three desiderata: (i) every
minority class is held by exactly two clients, creating co-trained
drifting pairs that the proximal anchor must reconcile; (ii) the
majority class melanocytic nevi is held by all seven clients,
providing a within-experiment per-class control where the proximal
anchor cannot help; (iii) per-client sample counts are within a
$1.65\times$ range. Three of seven clients co-train on
\emph{actinic keratoses + basal cell carcinoma}, two on
\emph{benign keratosis + dermatofibroma}, two on \emph{melanoma +
vascular lesions}; one client holds only nevi.

\begin{figure}[!htbp]
\centering
\includegraphics[width=\textwidth]{F5_partition_distributions.pdf}
\caption{Class distribution of the 7{,}007 DermaMNIST training
samples across the seven simulated clients under each of the four
partition designs. Cell colour shows the per-(client, class) sample
count on a $\log(1+\mathrm{count})$ scale; integer counts are
annotated. Per-client total $n$ is shown in the row label.}
\label{fig:partition-distributions}
\end{figure}

\subsection{Primary results: engineered partition}
\label{sec:headline-result}

The headline experiment compares FedAvg and FedProx on the
engineered \texttt{balanced\_paired\_7\_clients} partition through
the pure-PyTorch sequential reference loop. The protocol uses
ten paired seeds with identical partition, identical initial
weights, identical per-(round, client) dataloader RNG, and identical
optimiser state at $\mu = 0$ across the two algorithms. Within-pair
differences therefore reflect only the algorithmic difference (the
gated proximal term), not random-state divergence. Hyperparameters
are fixed at $E = 20$ local epochs, $R = 150$ communication rounds,
learning rate $0.01$, momentum $0.9$, $\mu = 0.01$, and full
participation $C = 1.0$. The proximal coefficient $\mu = 0.01$ was
chosen a~priori following \citet{li2020fedprox}; sparse exploratory
checks at $\mu \in \{0.001, 0.1\}$ on a subset of seeds are
consistent with $\mu = 0.01$ being near a local optimum.

\paragraph{Headline.}
Table~\ref{tab:headline-summary} reports test macro-F1 means $\pm$
standard deviations across the ten paired seeds. FedProx outperforms
FedAvg with mean macro-F1 of $0.508 \pm 0.014$ vs $0.481 \pm 0.025$,
a within-pair improvement of $+0.0267 \pm 0.0349$ macro-F1. FedProx
wins on 9 of 10 paired seeds; the lone losing seed yields
$\Delta = -0.020$, and the largest winning seed yields
$\Delta = +0.095$. FedProx's across-seed standard deviation
($0.014$) is roughly half that of FedAvg's ($0.025$), indicating
that FedProx produces more reproducible global models under non-IID
conditions.

% T01 -- Headline summary (field-standard format, auto-generated)
\begin{table}[!htbp]
\centering
\caption{Headline test macro-F1 on the engineered \texttt{balanced\_paired\_7\_clients} partition under the pure-PyTorch reference loop, reported as mean $\pm$ standard deviation across $n=10$ paired seeds. \textbf{Bolding convention used throughout the thesis tables:} a cell is bolded only when the corresponding paired Wilcoxon contrast is significant at $\alpha = 0.05$ (per-class tables: Holm-corrected over the family of seven classes). Non-bold rows are within seed-noise.}
\label{tab:headline-summary}
\small
\begin{tabular}{lcc}
\toprule
& FedAvg & FedProx \\
\midrule
Test macro-F1                 & $0.4814 \pm 0.0254$ & \textbf{$0.5081 \pm 0.0140$} \\
Mean within-pair $\Delta$     & \multicolumn{2}{c}{$+0.0267 \pm 0.0349$ \quad (paired Wilcoxon $p = 0.020$, Walsh 95\% CI $[+0.002,\,+0.053]$)} \\
FedProx wins                  & \multicolumn{2}{c}{9/10 paired seeds} \\
\bottomrule
\end{tabular}
\end{table}


\begin{figure}[!htbp]
\centering
\includegraphics[width=\textwidth]{F1_headline_convergence.pdf}
\caption{Headline convergence on the engineered partition
(pure-PyTorch, $n = 10$ paired seeds, mean $\pm$ SEM band).
\textbf{Left:} validation macro-F1 over the 150 communication
rounds; FedProx (dark slate blue) maintains a higher trajectory
than FedAvg (mint green) from round $\sim 10$ onwards.
\textbf{Centre:} validation cross-entropy loss; both algorithms
exhibit overfitting after round $\sim 30$ but FedProx's loss
remains $0.10$--$0.18$ below FedAvg's throughout late training.
\textbf{Right:} test macro-F1 at the best-validation checkpoint
(solid) and at the final round (hatched), with the centralised
pure-PyTorch reference (deep navy) shown for context. FedProx
recovers $\sim 91\%$ of centralised performance; FedAvg
$\sim 86\%$.}
\label{fig:headline-convergence}
\end{figure}

\paragraph{Per-class breakdown.}
Table~\ref{tab:per-class} and Figure~\ref{fig:per-class-convergence}
decompose the headline result by class. The FedProx advantage is
strongly concentrated on the \emph{melanoma} class: FedProx
$0.301 \pm 0.026$ vs FedAvg $0.187 \pm 0.101$, an improvement of
$+0.114$ macro-F1 with FedProx winning on all 10 paired seeds. The
melanoma improvement alone accounts for roughly four times the
overall macro-F1 gain when averaged across the seven classes. The
next-largest improvement is on actinic keratoses
($+0.067$), followed by smaller positive differences on the
remaining minority classes.

The majority class \emph{melanocytic nevi} shows essentially zero
difference ($+0.002$ macro-F1) and the smallest across-seed standard
deviation of any class ($\pm 0.003$ for FedAvg, $\pm 0.004$ for
FedProx). This is consistent with the partition-design mechanism
prediction: melanocytic nevi is held by every client, so no
inter-client drift exists on this class for the proximal anchor to
reconcile, and FedProx has no opportunity to differentiate from
FedAvg on it. The two algorithms produce essentially identical
melanocytic-nevi trajectories
(Figure~\ref{fig:per-class-convergence}, centre row, middle), in
contrast to the visible separation on minority classes held by
client pairs.

% T02 -- Per-class test macro-F1 (field-standard, auto-generated)
\begin{table}[!htbp]
\centering
\caption{Per-class test F1 on the engineered partition under the pure-PyTorch reference loop, reported as mean $\pm$ standard deviation across $n=10$ paired seeds. Class prevalence in the 7\,007-sample training set is shown in parentheses. The winning algorithm per class is shown in bold; $\Delta$ is the mean within-pair difference (FedProx $-$ FedAvg).}
\label{tab:per-class}
\small
\begin{tabular}{lccc}
\toprule
Class (prevalence) & FedAvg & FedProx & $\Delta$ \\
\midrule
Actinic keratoses (3.27\%) & $0.404 \pm 0.111$ & \textbf{$0.471 \pm 0.040$} & $+0.067$ \\
Basal cell carcinoma (5.13\%) & \textbf{$0.475 \pm 0.046$} & $0.473 \pm 0.027$ & $-0.001$ \\
Benign keratosis (10.97\%) & $0.440 \pm 0.033$ & \textbf{$0.444 \pm 0.041$} & $+0.004$ \\
Dermatofibroma (1.15\%) & $0.283 \pm 0.141$ & \textbf{$0.296 \pm 0.070$} & $+0.013$ \\
Melanoma (11.11\%) & $0.187 \pm 0.101$ & \textbf{$0.301 \pm 0.026$} & $+0.114$ \\
Melanocytic nevi (67.05\%) & $0.884 \pm 0.003$ & \textbf{$0.886 \pm 0.004$} & $+0.002$ \\
Vascular lesions (1.41\%) & \textbf{$0.698 \pm 0.078$} & $0.686 \pm 0.043$ & $-0.012$ \\
\bottomrule
\end{tabular}
\end{table}


\begin{figure}[!htbp]
\centering
\includegraphics[width=\textwidth]{F2_per_class_convergence.pdf}
\caption{Per-class validation F1 over rounds (pure-PyTorch headline,
$n = 10$ paired seeds, mean $\pm$ SEM). The melanoma panel shows
the largest persistent gap between FedProx and FedAvg, while the
majority class melanocytic nevi (centre, middle) shows essentially
identical trajectories --- consistent with the partition design's
mechanism prediction that the proximal anchor can only act where
inter-client drift exists.}
\label{fig:per-class-convergence}
\end{figure}

\subsubsection{Direct mechanism evidence: per-client update norms}
\label{sec:update-norms}

The above results infer the proximal-anchor mechanism from
\emph{outcomes} (test F1 patterns consistent with drift control).
We additionally measure the mechanism directly by logging each
client's per-round update norm $\lVert w_i^{t+1} - w^t \rVert_2$ on
the Flower runtime headline sweep. The proximal anchor is
mathematically expected to reduce these norms because the penalty
term $(\mu/2)\lVert w - w^t\rVert^2$ acts as a proximal regulariser
on the parameter increment.

Figure~\ref{fig:update-norms} confirms this prediction directly.
Across all 150 communication rounds, FedAvg's mean per-client
update norm ($1.52$) is approximately $1.48\times$ that of
FedProx's ($1.03$). The gap is present from round~1 and persists
throughout training; the relative gap widens slightly in late
training as both algorithms approach a plateau but the proximal
constraint becomes a larger fraction of the local gradient. A
related outcome-level signature is the validation-trajectory
volatility: the mean absolute round-to-round change in validation
macro-F1 is approximately $1.84\times$ larger for FedAvg than
FedProx, with the ratio exceeding one on every paired seed. The
proximal anchor produces smaller per-round parameter updates which
translate into smoother per-round validation trajectories ---
exactly the mechanism's empirical signature.

\begin{figure}[!htbp]
\centering
\includegraphics[width=\textwidth]{F10_update_norms.pdf}
\caption{Mean per-client per-round update norm
$\lVert w_i^{t+1} - w^t\rVert_2$ across the 150 communication
rounds, averaged over clients within each round and then over
paired seeds (Flower runtime, engineered partition, mean $\pm$ SEM
band). FedProx produces systematically smaller per-client updates
from round~1 onwards, with an across-all-rounds mean ratio of
approximately $1.48\times$.}
\label{fig:update-norms}
\end{figure}

\subsubsection{Peak-vs-final-round behaviour}
\label{sec:peak-vs-final}

Table~\ref{tab:best-vs-final} compares the peak validation macro-F1
reached by each algorithm with the final-round validation macro-F1.
Both algorithms exhibit a similar post-peak drop of approximately
$0.04$--$0.05$ macro-F1, with FedAvg's drop ($0.046$) marginally
larger than FedProx's ($0.043$). FedProx's higher trajectory
throughout late training (Figure~\ref{fig:headline-convergence})
is therefore best understood as a parallel upward shift of the
entire validation curve rather than a reduced overfitting tendency
in the final third of training. The mechanism producing the
headline advantage operates throughout the training horizon, not
only in the post-peak regime.

% T11 -- Best-vs-final-round (field-standard, auto-generated)
\begin{table}[!htbp]
\centering
\caption{Peak vs final-round validation macro-F1 across the 150 communication rounds (pure-PyTorch headline, $n = 10$ paired seeds, engineered partition). The post-peak drop quantifies how much each algorithm overfits between its best and final round.}
\label{tab:best-vs-final}
\small
\begin{tabular}{lccc}
\toprule
Algorithm & Peak val F1 & Final val F1 & Post-peak drop \\
\midrule
FedAvg  & $0.5188 \pm 0.0174$ & $0.4729 \pm 0.0335$ & $+0.0459 \pm 0.0357$ \\
FedProx & \textbf{$0.5540 \pm 0.0271$} & \textbf{$0.5114 \pm 0.0415$} & $+0.0425 \pm 0.0339$ \\
\bottomrule
\end{tabular}
\end{table}


\subsection{Robustness across partition designs}
\label{sec:partition-robustness}

Table~\ref{tab:partition-robustness} and
Figure~\ref{fig:partition-convergence} report the FedProx-vs-FedAvg
comparison on the four partition designs, all on the Flower
simulation runtime with the same ten paired seeds.

% T04 -- Partition robustness (field-standard, auto-generated)
\begin{table}[!htbp]
\centering
\caption{Robustness across four DermaMNIST partition designs, all on the Flower simulation runtime with the same ten paired seeds. The IID partition is the theoretical-null condition; the specialist row is the pre-registered specialist counterfactual described in \S\ref{sec:hyp-specialist} (same client sample counts $n_i$ as the engineered partition, every minority class assigned to exactly one client rather than two; included to probe whether the pairing structure --- not just label skew --- drives the FedProx effect); Dirichlet and balanced-paired are descriptive robustness rows. Values are mean $\pm$ SD across seeds; $\Delta$ is the within-pair FedProx $-$ FedAvg difference; $p$ is the two-sided paired Wilcoxon $p$-value. \textbf{No row is bolded} because none of the four per-partition contrasts reaches $\alpha = 0.05$ at $n = 10$ paired seeds. The IID row's slightly negative descriptive $\Delta$ matches the theoretical-null prediction; the specialist row's magnitude relative to balanced-paired is inferentially indeterminate at this sample size (\S\ref{sec:partition-robustness}).}
\label{tab:partition-robustness}
\small
\begin{tabular}{lccccc}
\toprule
Partition & FedAvg & FedProx & $\Delta$ & $p$ (Wilcoxon) & Wins \\
\midrule
IID (uniform random) & $0.5853 \pm 0.0205$ & $0.5785 \pm 0.0204$ & $-0.0068$ & $> 0.05$ & 4/10 \\
Dirichlet $\alpha = 0.1$ & $0.4796 \pm 0.0355$ & $0.4875 \pm 0.0341$ & $+0.0079$ & $> 0.05$ & 6/10 \\
Specialist (1-of-7) & $0.4365 \pm 0.0236$ & $0.4446 \pm 0.0172$ & $+0.0081$ & $> 0.05$ & 7/10 \\
Balanced paired (engineered) & $0.4964 \pm 0.0150$ & $0.5033 \pm 0.0206$ & $+0.0069$ & $0.432$ & 7/10 \\
\bottomrule
\end{tabular}
\end{table}


\paragraph{IID partition.}
Under uniform random partitioning, FedProx slightly underperforms
FedAvg ($-0.007$ mean within-pair difference, FedAvg wins 6/10
paired seeds). The two algorithms' validation trajectories are
visually indistinguishable from round $\sim 20$ onwards
(Figure~\ref{fig:partition-convergence}, top left). This is the
predicted behaviour of the proximal-anchor mechanism: when local
objectives are i.i.d.\ samples of the same distribution, no
inter-client drift exists and the proximal term acts only as a
small regularisation cost on the local optimisation. The slight
negative sign of $\Delta$ under IID is mechanistically consistent
with this regularisation overhead.

\paragraph{Non-IID partitions.}
Across the three non-IID partition designs (Dirichlet
$\alpha = 0.1$, specialist, balanced paired) FedProx descriptively
outperforms FedAvg with within-pair differences of $+0.008$,
$+0.008$, and $+0.007$ respectively. The differences are small in
absolute magnitude and the across-seed standard deviations are
comparable in size to the differences themselves, so the magnitude
is not reliably separable from seed-level variation at $n = 10$.
However, the direction is consistent across all three designs and
the sign of $\Delta$ is opposite to the IID case --- the proximal
anchor produces a positive effect whenever class skew is present,
regardless of the specific form of that skew.

\paragraph{Pairing counterfactual.}
The specialist partition was constructed to isolate the
\emph{pairing} feature of the engineered partition from generic
class specialisation. Each minority class is held by exactly one
client under specialist (no within-pair co-training), versus exactly
two clients under engineered. The descriptive comparison shows the
specialist $\Delta$ ($+0.008$) is similar in magnitude to the
engineered $\Delta$ on the same runtime ($+0.007$), suggesting that
generic singleton class assignment is sufficient to produce the
small FedProx benefit visible on Flower at this sample size. The
pairing feature does not appear to further amplify the effect ---
though this conclusion is constrained by the small absolute
magnitudes and the across-seed spread.

\begin{figure}[!htbp]
\centering
\includegraphics[width=\textwidth]{F4_partition_convergence.pdf}
\caption{Partition robustness on the Flower runtime ($n = 10$
paired seeds, mean $\pm$ SEM). Under IID (top left) the two
algorithms are visually indistinguishable from round $\sim 20$
onwards. Under Dirichlet $\alpha = 0.1$, specialist, and
engineered balanced-paired partitions, FedProx maintains a modest
visible advantage of order $0.01$--$0.02$ macro-F1 throughout
training.}
\label{fig:partition-convergence}
\end{figure}

\subsection{Cross-runtime sensitivity}
\label{sec:cross-runtime}

The headline experiment is replicated on the Flower simulation
runtime using the same engineered partition and the same ten paired
seeds. The pure-PyTorch reference loop is a single-process
sequential implementation; the Flower runtime uses Ray actors with
inter-process parameter serialisation. Both implement bit-equivalent
local objectives at $\mu = 0$ and use identical hyperparameters.

Under pure-PyTorch the within-pair difference is
$\Delta = +0.027 \pm 0.035$ (FedProx wins 9/10 seeds); under Flower
it is $\Delta = +0.007 \pm 0.024$ (FedProx wins 7/10).
Table~\ref{tab:cross-runtime} reports both side by side. The
direction is consistent across runtimes (positive $\Delta$, FedProx
wins majority of paired seeds), but the across-seed mean magnitude
is approximately $4\times$ larger under pure-PyTorch. The per-seed
gap (pure-PyTorch $\Delta$ minus Flower $\Delta$) is not consistent
across the ten seeds: five seeds show pure-PyTorch $>$ Flower and
five show Flower $\geq$ pure-PyTorch. The cross-runtime difference
in magnitude therefore appears to be a high-variance noise effect
attributable to floating-point non-associativity in Flower's Ray-
actor aggregation, rather than a systematic per-seed shift.

% T03 -- Cross-runtime comparison (field-standard, auto-generated)
\begin{table}[!htbp]
\centering
\caption{Cross-runtime replication on the engineered \texttt{balanced\_paired\_7\_clients} partition. Both runtimes use the same ten paired seeds and identical hyperparameters; the only difference is the implementation (pure-PyTorch sequential reference loop vs Flower simulation with Ray actors). Values are mean $\pm$ SD across seeds. $p$ is the within-runtime two-sided paired Wilcoxon $p$-value. \textbf{Bold} marks within-runtime contrasts significant at $\alpha = 0.05$; the Flower-runtime contrast is descriptively positive but not significant.}
\label{tab:cross-runtime}
\small
\begin{tabular}{lccccc}
\toprule
Runtime & FedAvg & FedProx & $\Delta$ & $p$ (Wilcoxon) & Wins \\
\midrule
Pure-PyTorch (reference) & $0.4814 \pm 0.0254$ & \textbf{$0.5081 \pm 0.0140$} & \textbf{$+0.0267$} & \textbf{$0.020$} & 9/10 \\
Flower simulation & $0.4964 \pm 0.0150$ & $0.5033 \pm 0.0206$ & $+0.0069$ & $0.432$ & 7/10 \\
\bottomrule
\end{tabular}
\end{table}


\begin{figure}[!htbp]
\centering
\includegraphics[width=\textwidth]{F3_cross_runtime_convergence.pdf}
\caption{Cross-runtime convergence on the engineered partition
($n = 10$ paired seeds, mean $\pm$ SEM). Both the pure-PyTorch
sequential loop (left) and the Flower simulation runtime (right)
show FedProx maintaining a higher trajectory than FedAvg throughout
training. The across-seed mean within-pair gap is larger under
pure-PyTorch ($+0.027$) than under Flower ($+0.007$).}
\label{fig:cross-runtime-convergence}
\end{figure}

\subsection{Federation tax}
\label{sec:federation-tax}

To anchor the absolute magnitudes of the federated results, we
trained a centralised non-federated reference on the pooled
7{,}007-sample training set under the same architecture, optimiser,
and evaluation protocol. Table~\ref{tab:federation-tax} reports
test macro-F1 for the centralised reference and for the four
federated configurations at the matched ten seeds.

% T05 -- Federation tax (field-standard, auto-generated)
\begin{table}[!htbp]
\centering
\caption{Federation tax: comparison of centralised training (pooled non-federated, upper-bound reference) against the four federated configurations on the engineered \texttt{balanced\_paired\_7\_clients} partition. The recovery column expresses each FL configuration's mean as a fraction of the centralised mean.}
\label{tab:federation-tax}
\small
\begin{tabular}{lcc}
\toprule
Configuration & Test macro-F1 (mean $\pm$ SD) & Recovery vs centralised \\
\midrule
Centralised (no FL) & \textbf{$0.5600 \pm 0.0240$} & --- \\
FedAvg (pure-PyTorch) & $0.4814 \pm 0.0254$ & 86.0\% \\
FedProx (pure-PyTorch) & $0.5081 \pm 0.0140$ & 90.7\% \\
FedAvg (Flower) & $0.4964 \pm 0.0150$ & 88.6\% \\
FedProx (Flower) & $0.5033 \pm 0.0206$ & 89.9\% \\
\bottomrule
\end{tabular}
\end{table}


The centralised mean is $0.560 \pm 0.024$, the upper bound for any
FL configuration on this task. Pure-PyTorch FedAvg recovers
$86.0\%$ of centralised performance; pure-PyTorch FedProx recovers
$90.7\%$; Flower FedAvg $88.6\%$; Flower FedProx $89.9\%$. The
within-runtime improvement from FedAvg to FedProx (approximately
$5$ percentage points under pure-PyTorch, $1$--$2$ under Flower) is
small compared to the federation tax itself
($9$--$14$ percentage points). The headline FedProx advantage is
therefore best framed not as ``federated training matches
centralised'' but as ``federated FedProx closes a modest portion of
the gap that federated FedAvg leaves open.''

\subsection{Summary}
\label{sec:headline-summary}

The pre-registered primary analysis on the engineered partition
under pure-PyTorch yields a $+0.027$ macro-F1 within-pair improvement
for FedProx over FedAvg, with FedProx winning 9 of 10 paired seeds.
The improvement is concentrated on the melanoma class
($\Delta = +0.114$, 10/10 paired wins), the clinically most
important diagnosis, while the majority class melanocytic nevi
shows essentially no difference --- consistent with the
proximal-anchor mechanism's prediction that the regulariser is
inert when every client holds the class. Direct measurement of
per-client update norms confirms FedProx produces approximately
$1.48\times$ smaller per-round parameter updates than FedAvg,
providing empirical mechanism evidence for the drift-control
interpretation. The improvement direction is preserved across the
three additional non-IID partition designs (Dirichlet, specialist,
engineered on Flower), all with positive $\Delta$ of order
$+0.007$--$+0.008$; the IID partition's small negative $\Delta$
($-0.007$) is consistent with the mechanism's prediction that the
proximal anchor cannot help when no inter-client drift exists. The
cross-runtime replication on Flower preserves the direction of
effect but produces a smaller magnitude, reflecting Ray-actor
non-determinism in aggregation.

\paragraph{Scope and limitations.}
The headline partition is engineered to maximise the proximal-anchor
mechanism's effect; its external validity to real multi-institutional
cohorts is unestablished and explicit. Effect magnitudes are smaller
than reported in the original \citet{li2020fedprox} benchmarks
(Synthetic, FEMNIST, Shakespeare, Sentiment140), consistent with
DermaMNIST being a smaller and more challenging task than those
benchmarks. The across-seed standard deviations are comparable in
size to the mean within-pair differences on the Flower runtime,
which limits the strength of per-partition conclusions at the chosen
sample size of $n = 10$. The melanoma per-class result is the
single most robust per-class finding; per-class differences for the
remaining minority classes are small relative to seed-level
variation and should be read as suggestive rather than definitive.
The DermaMNIST partitions used here are synthetic simulations of
non-i.i.d.\ structure rather than empirical multi-institutional
data; whether the magnitudes observed here transfer to real
multi-site federated deployments is identified as future work.

% ==============================================================================

\section{Statistical Heterogeneity}
\label{sec:results-statistical-het}

Statistical heterogeneity refers to differences in the data
distributions held by the participating clients. In medical-imaging
federated learning this is the dominant source of non-i.i.d.\
behaviour: each institution sees a different case mix, so per-client
class priors $p_i(y)$ diverge from the global prior $p(y)$. The
proximal term in FedProx is theoretically motivated by this regime
\citep{li2020fedprox}: when local objectives $F_i$ disagree at the
broadcast iterate $w^t$, the proximal anchor
$(\mu/2)\lVert w - w^t \rVert^2$ damps client drift and keeps the
aggregated update closer to a consensus that minimises the federated
objective $F$.

This section reports the FedProx-vs-FedAvg comparison under four
partition designs of DermaMNIST, each chosen to expose a different
facet of label heterogeneity. We follow the reporting convention of
the federated-learning literature \citep{li2020fedprox, mcmahan2017fedavg,
wang2020fednova, hsu2019dirichlet}: descriptive comparison of means
(mean $\pm$ SD across paired seeds), visual curve comparisons with
shaded uncertainty bands, and bolded best results in tables, without
formal hypothesis testing. The headline
\texttt{balanced\_paired\_7\_clients} partition is the primary
analysis, run on the pure-PyTorch reference loop. The three other
partitions (IID, Dirichlet $\alpha = 0.1$, specialist) are evaluated
on the Flower simulation runtime as robustness checks; the engineered
partition is also replicated on Flower for the cross-runtime
sensitivity check (Section~\ref{sec:cross-runtime}).

\subsection{Experimental design and partition motivation}
\label{sec:partition-designs}

Each partition design corresponds to a specific question about the
FedProx mechanism. The class-distribution visualisations are shown
in Figure~\ref{fig:partition-distributions}.

\paragraph{IID partition (Figure~\ref{fig:partition-distributions}a) ---
falsification check.}
Each of the seven clients receives a uniform random sample of $\sim
1{,}001$ training images drawn i.i.d.\ from the global DermaMNIST
pool. The per-client class distribution mirrors the global
distribution in expectation, with the majority class (melanocytic
nevi, 67\%) proportionally represented in every client. Under this
partition the local objectives $F_i$ are i.i.d.\ samples of the same
distribution and their gradients at any $w^t$ are unbiased estimators
of $\nabla F(w^t)$; the proximal anchor has no client drift to
constrain and FedProx is predicted to provide no advantage over
FedAvg ($\Delta \approx 0$). A positive $\Delta$ under IID would
contradict the drift-control mechanism interpretation, so this arm
functions as a null-condition probe.

\paragraph{Dirichlet $\alpha = 0.1$ partition
(Figure~\ref{fig:partition-distributions}b) --- literature-standard
non-i.i.d.}
The Dirichlet-partitioning protocol of \citet{hsu2019dirichlet} is
the de-facto standard for synthetic non-i.i.d.\ benchmarking in
federated learning. For each class $c$ a vector
$\bm{q}_c \sim \mathrm{Dir}_7(\alpha \bm{1})$ is drawn and the
class-$c$ samples are allocated to the seven clients in proportions
$q_{c,1}, \dots, q_{c,7}$. The concentration $\alpha = 0.1$ produces
highly skewed priors: in the realisation shown, client~6 receives
3{,}721 samples (mostly nevi) while client~4 receives only~91. This
combines label skew with severe quantity skew --- the two axes
described in \citet{kairouz2021advances} --- and is the regime in
which most non-i.i.d.\ FL benchmarks report FedAvg degradation.

\paragraph{Specialist partition (Figure~\ref{fig:partition-distributions}c)
--- pairing-lever counterfactual.}
Every minority class is assigned to exactly one client; the majority
class (melanocytic nevi) is split across all seven clients with
unequal shares chosen so per-client sample counts $n_i$ match the
engineered balanced-paired partition exactly. The specialist design
isolates the \emph{paired-co-training} feature from a generic
label-skew effect: under the engineered partition every minority
class is held by two clients (which must drift apart during local
SGD and be reconciled at aggregation), whereas under specialist each
minority class is held by exactly one client (no inter-client drift
exists for that class). The contrast tests whether the pairing
structure or the more generic specialisation is the mechanism's
target.

\paragraph{Balanced paired partition (Figure~\ref{fig:partition-distributions}d)
--- engineered headline.}
The headline partition satisfies three desiderata: (i) every
minority class is held by exactly two clients, creating co-trained
drifting pairs that the proximal anchor must reconcile; (ii) the
majority class melanocytic nevi is held by all seven clients,
providing a within-experiment per-class control where the proximal
anchor cannot help; (iii) per-client sample counts are within a
$1.65\times$ range. Three of seven clients co-train on
\emph{actinic keratoses + basal cell carcinoma}, two on
\emph{benign keratosis + dermatofibroma}, two on \emph{melanoma +
vascular lesions}; one client holds only nevi.

\begin{figure}[!htbp]
\centering
\includegraphics[width=\textwidth]{F5_partition_distributions.pdf}
\caption{Class distribution of the 7{,}007 DermaMNIST training
samples across the seven simulated clients under each of the four
partition designs. Cell colour shows the per-(client, class) sample
count on a $\log(1+\mathrm{count})$ scale; integer counts are
annotated. Per-client total $n$ is shown in the row label.}
\label{fig:partition-distributions}
\end{figure}

\subsection{Primary results: engineered partition}
\label{sec:headline-result}

The headline experiment compares FedAvg and FedProx on the
engineered \texttt{balanced\_paired\_7\_clients} partition through
the pure-PyTorch sequential reference loop. The protocol uses
ten paired seeds with identical partition, identical initial
weights, identical per-(round, client) dataloader RNG, and identical
optimiser state at $\mu = 0$ across the two algorithms. Within-pair
differences therefore reflect only the algorithmic difference (the
gated proximal term), not random-state divergence. Hyperparameters
are fixed at $E = 20$ local epochs, $R = 150$ communication rounds,
learning rate $0.01$, momentum $0.9$, $\mu = 0.01$, and full
participation $C = 1.0$. The proximal coefficient $\mu = 0.01$ was
chosen a~priori following \citet{li2020fedprox}; sparse exploratory
checks at $\mu \in \{0.001, 0.1\}$ on a subset of seeds are
consistent with $\mu = 0.01$ being near a local optimum.

\paragraph{Headline.}
Table~\ref{tab:headline-summary} reports test macro-F1 means $\pm$
standard deviations across the ten paired seeds. FedProx outperforms
FedAvg with mean macro-F1 of $0.508 \pm 0.014$ vs $0.481 \pm 0.025$,
a within-pair improvement of $+0.0267 \pm 0.0349$ macro-F1. FedProx
wins on 9 of 10 paired seeds; the lone losing seed yields
$\Delta = -0.020$, and the largest winning seed yields
$\Delta = +0.095$. FedProx's across-seed standard deviation
($0.014$) is roughly half that of FedAvg's ($0.025$), indicating
that FedProx produces more reproducible global models under non-IID
conditions.

% T01 -- Headline summary (field-standard format, auto-generated)
\begin{table}[!htbp]
\centering
\caption{Headline test macro-F1 on the engineered \texttt{balanced\_paired\_7\_clients} partition under the pure-PyTorch reference loop, reported as mean $\pm$ standard deviation across $n=10$ paired seeds. \textbf{Bolding convention used throughout the thesis tables:} a cell is bolded only when the corresponding paired Wilcoxon contrast is significant at $\alpha = 0.05$ (per-class tables: Holm-corrected over the family of seven classes). Non-bold rows are within seed-noise.}
\label{tab:headline-summary}
\small
\begin{tabular}{lcc}
\toprule
& FedAvg & FedProx \\
\midrule
Test macro-F1                 & $0.4814 \pm 0.0254$ & \textbf{$0.5081 \pm 0.0140$} \\
Mean within-pair $\Delta$     & \multicolumn{2}{c}{$+0.0267 \pm 0.0349$ \quad (paired Wilcoxon $p = 0.020$, Walsh 95\% CI $[+0.002,\,+0.053]$)} \\
FedProx wins                  & \multicolumn{2}{c}{9/10 paired seeds} \\
\bottomrule
\end{tabular}
\end{table}


\begin{figure}[!htbp]
\centering
\includegraphics[width=\textwidth]{F1_headline_convergence.pdf}
\caption{Headline convergence on the engineered partition
(pure-PyTorch, $n = 10$ paired seeds, mean $\pm$ SEM band).
\textbf{Left:} validation macro-F1 over the 150 communication
rounds; FedProx (dark slate blue) maintains a higher trajectory
than FedAvg (mint green) from round $\sim 10$ onwards.
\textbf{Centre:} validation cross-entropy loss; both algorithms
exhibit overfitting after round $\sim 30$ but FedProx's loss
remains $0.10$--$0.18$ below FedAvg's throughout late training.
\textbf{Right:} test macro-F1 at the best-validation checkpoint
(solid) and at the final round (hatched), with the centralised
pure-PyTorch reference (deep navy) shown for context. FedProx
recovers $\sim 91\%$ of centralised performance; FedAvg
$\sim 86\%$.}
\label{fig:headline-convergence}
\end{figure}

\paragraph{Per-class breakdown.}
Table~\ref{tab:per-class} and Figure~\ref{fig:per-class-convergence}
decompose the headline result by class. The FedProx advantage is
strongly concentrated on the \emph{melanoma} class: FedProx
$0.301 \pm 0.026$ vs FedAvg $0.187 \pm 0.101$, an improvement of
$+0.114$ macro-F1 with FedProx winning on all 10 paired seeds. The
melanoma improvement alone accounts for roughly four times the
overall macro-F1 gain when averaged across the seven classes. The
next-largest improvement is on actinic keratoses
($+0.067$), followed by smaller positive differences on the
remaining minority classes.

The majority class \emph{melanocytic nevi} shows essentially zero
difference ($+0.002$ macro-F1) and the smallest across-seed standard
deviation of any class ($\pm 0.003$ for FedAvg, $\pm 0.004$ for
FedProx). This is consistent with the partition-design mechanism
prediction: melanocytic nevi is held by every client, so no
inter-client drift exists on this class for the proximal anchor to
reconcile, and FedProx has no opportunity to differentiate from
FedAvg on it. The two algorithms produce essentially identical
melanocytic-nevi trajectories
(Figure~\ref{fig:per-class-convergence}, centre row, middle), in
contrast to the visible separation on minority classes held by
client pairs.

% T02 -- Per-class test macro-F1 (field-standard, auto-generated)
\begin{table}[!htbp]
\centering
\caption{Per-class test F1 on the engineered partition under the pure-PyTorch reference loop, reported as mean $\pm$ standard deviation across $n=10$ paired seeds. Class prevalence in the 7\,007-sample training set is shown in parentheses. The winning algorithm per class is shown in bold; $\Delta$ is the mean within-pair difference (FedProx $-$ FedAvg).}
\label{tab:per-class}
\small
\begin{tabular}{lccc}
\toprule
Class (prevalence) & FedAvg & FedProx & $\Delta$ \\
\midrule
Actinic keratoses (3.27\%) & $0.404 \pm 0.111$ & \textbf{$0.471 \pm 0.040$} & $+0.067$ \\
Basal cell carcinoma (5.13\%) & \textbf{$0.475 \pm 0.046$} & $0.473 \pm 0.027$ & $-0.001$ \\
Benign keratosis (10.97\%) & $0.440 \pm 0.033$ & \textbf{$0.444 \pm 0.041$} & $+0.004$ \\
Dermatofibroma (1.15\%) & $0.283 \pm 0.141$ & \textbf{$0.296 \pm 0.070$} & $+0.013$ \\
Melanoma (11.11\%) & $0.187 \pm 0.101$ & \textbf{$0.301 \pm 0.026$} & $+0.114$ \\
Melanocytic nevi (67.05\%) & $0.884 \pm 0.003$ & \textbf{$0.886 \pm 0.004$} & $+0.002$ \\
Vascular lesions (1.41\%) & \textbf{$0.698 \pm 0.078$} & $0.686 \pm 0.043$ & $-0.012$ \\
\bottomrule
\end{tabular}
\end{table}


\begin{figure}[!htbp]
\centering
\includegraphics[width=\textwidth]{F2_per_class_convergence.pdf}
\caption{Per-class validation F1 over rounds (pure-PyTorch headline,
$n = 10$ paired seeds, mean $\pm$ SEM). The melanoma panel shows
the largest persistent gap between FedProx and FedAvg, while the
majority class melanocytic nevi (centre, middle) shows essentially
identical trajectories --- consistent with the partition design's
mechanism prediction that the proximal anchor can only act where
inter-client drift exists.}
\label{fig:per-class-convergence}
\end{figure}

\subsubsection{Direct mechanism evidence: per-client update norms}
\label{sec:update-norms}

The above results infer the proximal-anchor mechanism from
\emph{outcomes} (test F1 patterns consistent with drift control).
We additionally measure the mechanism directly by logging each
client's per-round update norm $\lVert w_i^{t+1} - w^t \rVert_2$ on
the Flower runtime headline sweep. The proximal anchor is
mathematically expected to reduce these norms because the penalty
term $(\mu/2)\lVert w - w^t\rVert^2$ acts as a proximal regulariser
on the parameter increment.

Figure~\ref{fig:update-norms} confirms this prediction directly.
Across all 150 communication rounds, FedAvg's mean per-client
update norm ($1.52$) is approximately $1.48\times$ that of
FedProx's ($1.03$). The gap is present from round~1 and persists
throughout training; the relative gap widens slightly in late
training as both algorithms approach a plateau but the proximal
constraint becomes a larger fraction of the local gradient. A
related outcome-level signature is the validation-trajectory
volatility: the mean absolute round-to-round change in validation
macro-F1 is approximately $1.84\times$ larger for FedAvg than
FedProx, with the ratio exceeding one on every paired seed. The
proximal anchor produces smaller per-round parameter updates which
translate into smoother per-round validation trajectories ---
exactly the mechanism's empirical signature.

\begin{figure}[!htbp]
\centering
\includegraphics[width=\textwidth]{F10_update_norms.pdf}
\caption{Mean per-client per-round update norm
$\lVert w_i^{t+1} - w^t\rVert_2$ across the 150 communication
rounds, averaged over clients within each round and then over
paired seeds (Flower runtime, engineered partition, mean $\pm$ SEM
band). FedProx produces systematically smaller per-client updates
from round~1 onwards, with an across-all-rounds mean ratio of
approximately $1.48\times$.}
\label{fig:update-norms}
\end{figure}

\subsubsection{Peak-vs-final-round behaviour}
\label{sec:peak-vs-final}

Table~\ref{tab:best-vs-final} compares the peak validation macro-F1
reached by each algorithm with the final-round validation macro-F1.
Both algorithms exhibit a similar post-peak drop of approximately
$0.04$--$0.05$ macro-F1, with FedAvg's drop ($0.046$) marginally
larger than FedProx's ($0.043$). FedProx's higher trajectory
throughout late training (Figure~\ref{fig:headline-convergence})
is therefore best understood as a parallel upward shift of the
entire validation curve rather than a reduced overfitting tendency
in the final third of training. The mechanism producing the
headline advantage operates throughout the training horizon, not
only in the post-peak regime.

% T11 -- Best-vs-final-round (field-standard, auto-generated)
\begin{table}[!htbp]
\centering
\caption{Peak vs final-round validation macro-F1 across the 150 communication rounds (pure-PyTorch headline, $n = 10$ paired seeds, engineered partition). The post-peak drop quantifies how much each algorithm overfits between its best and final round.}
\label{tab:best-vs-final}
\small
\begin{tabular}{lccc}
\toprule
Algorithm & Peak val F1 & Final val F1 & Post-peak drop \\
\midrule
FedAvg  & $0.5188 \pm 0.0174$ & $0.4729 \pm 0.0335$ & $+0.0459 \pm 0.0357$ \\
FedProx & \textbf{$0.5540 \pm 0.0271$} & \textbf{$0.5114 \pm 0.0415$} & $+0.0425 \pm 0.0339$ \\
\bottomrule
\end{tabular}
\end{table}


\subsection{Robustness across partition designs}
\label{sec:partition-robustness}

Table~\ref{tab:partition-robustness} and
Figure~\ref{fig:partition-convergence} report the FedProx-vs-FedAvg
comparison on the four partition designs, all on the Flower
simulation runtime with the same ten paired seeds.

% T04 -- Partition robustness (field-standard, auto-generated)
\begin{table}[!htbp]
\centering
\caption{Robustness across four DermaMNIST partition designs, all on the Flower simulation runtime with the same ten paired seeds. The IID partition is the theoretical-null condition; the specialist row is the pre-registered specialist counterfactual described in \S\ref{sec:hyp-specialist} (same client sample counts $n_i$ as the engineered partition, every minority class assigned to exactly one client rather than two; included to probe whether the pairing structure --- not just label skew --- drives the FedProx effect); Dirichlet and balanced-paired are descriptive robustness rows. Values are mean $\pm$ SD across seeds; $\Delta$ is the within-pair FedProx $-$ FedAvg difference; $p$ is the two-sided paired Wilcoxon $p$-value. \textbf{No row is bolded} because none of the four per-partition contrasts reaches $\alpha = 0.05$ at $n = 10$ paired seeds. The IID row's slightly negative descriptive $\Delta$ matches the theoretical-null prediction; the specialist row's magnitude relative to balanced-paired is inferentially indeterminate at this sample size (\S\ref{sec:partition-robustness}).}
\label{tab:partition-robustness}
\small
\begin{tabular}{lccccc}
\toprule
Partition & FedAvg & FedProx & $\Delta$ & $p$ (Wilcoxon) & Wins \\
\midrule
IID (uniform random) & $0.5853 \pm 0.0205$ & $0.5785 \pm 0.0204$ & $-0.0068$ & $> 0.05$ & 4/10 \\
Dirichlet $\alpha = 0.1$ & $0.4796 \pm 0.0355$ & $0.4875 \pm 0.0341$ & $+0.0079$ & $> 0.05$ & 6/10 \\
Specialist (1-of-7) & $0.4365 \pm 0.0236$ & $0.4446 \pm 0.0172$ & $+0.0081$ & $> 0.05$ & 7/10 \\
Balanced paired (engineered) & $0.4964 \pm 0.0150$ & $0.5033 \pm 0.0206$ & $+0.0069$ & $0.432$ & 7/10 \\
\bottomrule
\end{tabular}
\end{table}


\paragraph{IID partition.}
Under uniform random partitioning, FedProx slightly underperforms
FedAvg ($-0.007$ mean within-pair difference, FedAvg wins 6/10
paired seeds). The two algorithms' validation trajectories are
visually indistinguishable from round $\sim 20$ onwards
(Figure~\ref{fig:partition-convergence}, top left). This is the
predicted behaviour of the proximal-anchor mechanism: when local
objectives are i.i.d.\ samples of the same distribution, no
inter-client drift exists and the proximal term acts only as a
small regularisation cost on the local optimisation. The slight
negative sign of $\Delta$ under IID is mechanistically consistent
with this regularisation overhead.

\paragraph{Non-IID partitions.}
Across the three non-IID partition designs (Dirichlet
$\alpha = 0.1$, specialist, balanced paired) FedProx descriptively
outperforms FedAvg with within-pair differences of $+0.008$,
$+0.008$, and $+0.007$ respectively. The differences are small in
absolute magnitude and the across-seed standard deviations are
comparable in size to the differences themselves, so the magnitude
is not reliably separable from seed-level variation at $n = 10$.
However, the direction is consistent across all three designs and
the sign of $\Delta$ is opposite to the IID case --- the proximal
anchor produces a positive effect whenever class skew is present,
regardless of the specific form of that skew.

\paragraph{Pairing counterfactual.}
The specialist partition was constructed to isolate the
\emph{pairing} feature of the engineered partition from generic
class specialisation. Each minority class is held by exactly one
client under specialist (no within-pair co-training), versus exactly
two clients under engineered. The descriptive comparison shows the
specialist $\Delta$ ($+0.008$) is similar in magnitude to the
engineered $\Delta$ on the same runtime ($+0.007$), suggesting that
generic singleton class assignment is sufficient to produce the
small FedProx benefit visible on Flower at this sample size. The
pairing feature does not appear to further amplify the effect ---
though this conclusion is constrained by the small absolute
magnitudes and the across-seed spread.

\begin{figure}[!htbp]
\centering
\includegraphics[width=\textwidth]{F4_partition_convergence.pdf}
\caption{Partition robustness on the Flower runtime ($n = 10$
paired seeds, mean $\pm$ SEM). Under IID (top left) the two
algorithms are visually indistinguishable from round $\sim 20$
onwards. Under Dirichlet $\alpha = 0.1$, specialist, and
engineered balanced-paired partitions, FedProx maintains a modest
visible advantage of order $0.01$--$0.02$ macro-F1 throughout
training.}
\label{fig:partition-convergence}
\end{figure}

\subsection{Cross-runtime sensitivity}
\label{sec:cross-runtime}

The headline experiment is replicated on the Flower simulation
runtime using the same engineered partition and the same ten paired
seeds. The pure-PyTorch reference loop is a single-process
sequential implementation; the Flower runtime uses Ray actors with
inter-process parameter serialisation. Both implement bit-equivalent
local objectives at $\mu = 0$ and use identical hyperparameters.

Under pure-PyTorch the within-pair difference is
$\Delta = +0.027 \pm 0.035$ (FedProx wins 9/10 seeds); under Flower
it is $\Delta = +0.007 \pm 0.024$ (FedProx wins 7/10).
Table~\ref{tab:cross-runtime} reports both side by side. The
direction is consistent across runtimes (positive $\Delta$, FedProx
wins majority of paired seeds), but the across-seed mean magnitude
is approximately $4\times$ larger under pure-PyTorch. The per-seed
gap (pure-PyTorch $\Delta$ minus Flower $\Delta$) is not consistent
across the ten seeds: five seeds show pure-PyTorch $>$ Flower and
five show Flower $\geq$ pure-PyTorch. The cross-runtime difference
in magnitude therefore appears to be a high-variance noise effect
attributable to floating-point non-associativity in Flower's Ray-
actor aggregation, rather than a systematic per-seed shift.

% T03 -- Cross-runtime comparison (field-standard, auto-generated)
\begin{table}[!htbp]
\centering
\caption{Cross-runtime replication on the engineered \texttt{balanced\_paired\_7\_clients} partition. Both runtimes use the same ten paired seeds and identical hyperparameters; the only difference is the implementation (pure-PyTorch sequential reference loop vs Flower simulation with Ray actors). Values are mean $\pm$ SD across seeds. $p$ is the within-runtime two-sided paired Wilcoxon $p$-value. \textbf{Bold} marks within-runtime contrasts significant at $\alpha = 0.05$; the Flower-runtime contrast is descriptively positive but not significant.}
\label{tab:cross-runtime}
\small
\begin{tabular}{lccccc}
\toprule
Runtime & FedAvg & FedProx & $\Delta$ & $p$ (Wilcoxon) & Wins \\
\midrule
Pure-PyTorch (reference) & $0.4814 \pm 0.0254$ & \textbf{$0.5081 \pm 0.0140$} & \textbf{$+0.0267$} & \textbf{$0.020$} & 9/10 \\
Flower simulation & $0.4964 \pm 0.0150$ & $0.5033 \pm 0.0206$ & $+0.0069$ & $0.432$ & 7/10 \\
\bottomrule
\end{tabular}
\end{table}


\begin{figure}[!htbp]
\centering
\includegraphics[width=\textwidth]{F3_cross_runtime_convergence.pdf}
\caption{Cross-runtime convergence on the engineered partition
($n = 10$ paired seeds, mean $\pm$ SEM). Both the pure-PyTorch
sequential loop (left) and the Flower simulation runtime (right)
show FedProx maintaining a higher trajectory than FedAvg throughout
training. The across-seed mean within-pair gap is larger under
pure-PyTorch ($+0.027$) than under Flower ($+0.007$).}
\label{fig:cross-runtime-convergence}
\end{figure}

\subsection{Federation tax}
\label{sec:federation-tax}

To anchor the absolute magnitudes of the federated results, we
trained a centralised non-federated reference on the pooled
7{,}007-sample training set under the same architecture, optimiser,
and evaluation protocol. Table~\ref{tab:federation-tax} reports
test macro-F1 for the centralised reference and for the four
federated configurations at the matched ten seeds.

% T05 -- Federation tax (field-standard, auto-generated)
\begin{table}[!htbp]
\centering
\caption{Federation tax: comparison of centralised training (pooled non-federated, upper-bound reference) against the four federated configurations on the engineered \texttt{balanced\_paired\_7\_clients} partition. The recovery column expresses each FL configuration's mean as a fraction of the centralised mean.}
\label{tab:federation-tax}
\small
\begin{tabular}{lcc}
\toprule
Configuration & Test macro-F1 (mean $\pm$ SD) & Recovery vs centralised \\
\midrule
Centralised (no FL) & \textbf{$0.5600 \pm 0.0240$} & --- \\
FedAvg (pure-PyTorch) & $0.4814 \pm 0.0254$ & 86.0\% \\
FedProx (pure-PyTorch) & $0.5081 \pm 0.0140$ & 90.7\% \\
FedAvg (Flower) & $0.4964 \pm 0.0150$ & 88.6\% \\
FedProx (Flower) & $0.5033 \pm 0.0206$ & 89.9\% \\
\bottomrule
\end{tabular}
\end{table}


The centralised mean is $0.560 \pm 0.024$, the upper bound for any
FL configuration on this task. Pure-PyTorch FedAvg recovers
$86.0\%$ of centralised performance; pure-PyTorch FedProx recovers
$90.7\%$; Flower FedAvg $88.6\%$; Flower FedProx $89.9\%$. The
within-runtime improvement from FedAvg to FedProx (approximately
$5$ percentage points under pure-PyTorch, $1$--$2$ under Flower) is
small compared to the federation tax itself
($9$--$14$ percentage points). The headline FedProx advantage is
therefore best framed not as ``federated training matches
centralised'' but as ``federated FedProx closes a modest portion of
the gap that federated FedAvg leaves open.''

\subsection{Summary}
\label{sec:headline-summary}

The pre-registered primary analysis on the engineered partition
under pure-PyTorch yields a $+0.027$ macro-F1 within-pair improvement
for FedProx over FedAvg, with FedProx winning 9 of 10 paired seeds.
The improvement is concentrated on the melanoma class
($\Delta = +0.114$, 10/10 paired wins), the clinically most
important diagnosis, while the majority class melanocytic nevi
shows essentially no difference --- consistent with the
proximal-anchor mechanism's prediction that the regulariser is
inert when every client holds the class. Direct measurement of
per-client update norms confirms FedProx produces approximately
$1.48\times$ smaller per-round parameter updates than FedAvg,
providing empirical mechanism evidence for the drift-control
interpretation. The improvement direction is preserved across the
three additional non-IID partition designs (Dirichlet, specialist,
engineered on Flower), all with positive $\Delta$ of order
$+0.007$--$+0.008$; the IID partition's small negative $\Delta$
($-0.007$) is consistent with the mechanism's prediction that the
proximal anchor cannot help when no inter-client drift exists. The
cross-runtime replication on Flower preserves the direction of
effect but produces a smaller magnitude, reflecting Ray-actor
non-determinism in aggregation.

\paragraph{Scope and limitations.}
The headline partition is engineered to maximise the proximal-anchor
mechanism's effect; its external validity to real multi-institutional
cohorts is unestablished and explicit. Effect magnitudes are smaller
than reported in the original \citet{li2020fedprox} benchmarks
(Synthetic, FEMNIST, Shakespeare, Sentiment140), consistent with
DermaMNIST being a smaller and more challenging task than those
benchmarks. The across-seed standard deviations are comparable in
size to the mean within-pair differences on the Flower runtime,
which limits the strength of per-partition conclusions at the chosen
sample size of $n = 10$. The melanoma per-class result is the
single most robust per-class finding; per-class differences for the
remaining minority classes are small relative to seed-level
variation and should be read as suggestive rather than definitive.
The DermaMNIST partitions used here are synthetic simulations of
non-i.i.d.\ structure rather than empirical multi-institutional
data; whether the magnitudes observed here transfer to real
multi-site federated deployments is identified as future work.
